\documentclass[../main.tex]{subfiles}
\begin{document}
%==========================================TIÊU ĐỀ TẬP========================================
\begin{table}[h]
  \begin{adjustwidth}{-1cm}{}
    \begin{tabular}{>{\centering\arraybackslash}p{6cm}|>{\raggedright\arraybackslash}p{14cm}}
      \multirow{5}{6cm}
      % Cột bên trái: ÉPISODE và số tập
      {\\[-40pt]
        \flaregothic\fontsize{20pt}{20pt}\selectfont \centering
        \color{eptype}HISTOIRE\color{black}\\
        \gangofthree\fontsize{72pt}{72pt}\selectfont
        \color{epnum}8
        \\[18pt]
        % \flaregothic\fontsize{14pt}{14pt}\selectfont
        % \color{epattr}BIENVENUE, \uline{\color{Banpire} Mel} !
        % \flaregothic\fontsize{18pt}{18pt}\selectfont
        % \color{epattr}ÉPISODE PERDU
        \color{black}
      }
       &
      % Dòng thứ nhất: tên tập tiếng Nhật
      {\fotskip\fontsize{18pt}{18pt}\selectfont \ruby{三}{み}\ruby{日}{っか}\ruby{目}{め}の\ruby{寺}{てら}\ruby{参}{まい}り}
      \\[6pt]
      % Dòng thứ hai: tên tập tiếng Anh
       & {\cafeteria\fontsize{18pt}{18pt}\selectfont Temple Visit on the Third Day!}             \\[4pt]
      % Dòng thứ ba: tên tập tiếng Việt
       & {\vnmsans\fontsize{18pt}{18pt}\selectfont Mùng Ba Tết \hll\ - Đi lễ chùa đầu năm!}      \\[8pt]
      % Dòng thứ tư: Nhân vật xuất hiện
       & {\brian{\fontsize{14pt}{14pt}\selectfont Track used: Lunar Zoo Year - Ultimate Battle}}
    \end{tabular}
  \end{adjustwidth}
\end{table}
%===========================================NỘI DUNG==========================================
\color{black}

Buổi chiều hôm đó mùng Ba Tết, sau một buổi sáng đầy nhộn nhịp với bài kiểm tra Toán bất ngờ.

Kuon, Miko, Subaru, Shion và Fubuki cùng bố cậu về quê để thăm
mộ gia đình nhà Hikawa. Bố cậu đi xe máy, còn bốn người họ đi bằng cửa thần kỳ.

Cả nhóm xuất hiện trước lối vào nghĩa trang gia đình Hikawa, không khí nơi đây mang
một vẻ yên bình và thanh tịnh, được bao quanh bởi những thửa ruộng lúa trải dài. Mặt trời
buổi chiều tỏa ánh sáng dịu nhẹ, khiến sương mù mỏng manh lượn lờ trên mặt ruộng như
một tấm lụa trắng mềm mại. Tiếng chim ríu rít vang vọng khắp không gian, hòa quyện
với hương đồng gió nội.

Đúng lúc cả nhóm vừa ổn định lại sau khi bước ra từ cửa thần kỳ, một tiếng xe máy nổ
máy vang lên. Đó là Ryujin, ông đeo mũ bảo hiểm, ngồi trên chiếc xe đã khá cũ nhưng
vẫn chạy êm như ngày đầu.

Ông mỉm cười, nhẹ nhàng dựng xe dưới gốc cây gần đó và bước tới. Ánh mắt ông lướt
qua từng người, vừa ngạc nhiên vừa hài lòng khi thấy họ đã có mặt đúng lúc.

``Mấy đứa cũng tới kịp lúc ghê nhỉ.''

Ở trước mặt, con đường dẫn vào nghĩa trang là một lối đi nhỏ lát đá, hai bên rợp bóng
những cây bạch đàn cao vút. Gió nhẹ thổi qua, làm lá cây xào xạc, tạo nên âm thanh như
những lời thì thầm từ quá khứ xa xôi. Xa xa, bóng dáng các ngôi mộ thấp thoáng, mỗi
tấm bia đều được chăm chút cẩn thận, gợi lên sự kính trọng và lòng thành kính của con
cháu dành cho tổ tiên.

Miko - với tư cách là một cô nàng đã từng làm trong ngôi đền - ngước nhìn xung quanh,
ánh mắt tò mò nhưng cũng đầy tôn kính. Subaru hơi cúi đầu, như muốn tỏ sự trang nghiêm
trong không gian linh thiêng này.

Kuon đứng gần bố, ánh mắt cậu đầy sự tĩnh lặng. Nơi đây không chỉ là nơi an nghỉ của tổ
tiên cậu mà còn là nơi lưu giữ những ký ức quý giá, những câu chuyện về gia đình được
truyền từ thế hệ này sang thế hệ khác.
Cả nhóm đứng đó trong giây lát, hít một hơi thật sâu để cảm nhận không khí trong lành
của làng quê, sẵn sàng bước vào để làm tròn trách nhiệm của mình trong ngày đầu xuân
ý nghĩa này.

\dots Chỉ riêng Fubuki là người vẫn còn đang tỏ vẻ không nghiêm túc, khi trước mắt cô là
một khu đất rất rộng đến bao la: ``Địa điểm hoàn hảo để chơi trốn tìm đây rồi!''

Miko nghe vậy liền rén ngang: ``Đừng. Miko sợ lắm.'' Cậu nhướng mày, khẽ nhắc nhẹ
nàng Cáo trắng tinh nghịch đó: ``Em biết ta đến đây không phải để chơi đùa chứ?''

Bố cậu đi vào trong, ngoái đầu lại nhắc: ``Vào đi các con.'' Năm người đồng loạt bước
theo sau ông, con đường mòn nhỏ phía trước rợp bóng cây, ánh sáng xuyên qua từng kẽ
lá tạo thành những vệt loang lổ trên mặt đất. Lá cây khẽ lay động trong gió, thi thoảng lại
nghe tiếng rụng nhẹ nhàng.

``Nghĩa trang ở tận bên trong sao?'' nàng Vịt ngạc nhiên, khi xung quanh họ là những rạp
cây bạch đàn, hai hàng ruộng lúa đã bắt đầu mơn mởn những chồi lúa non đầu tiên.

``Thấy ớn lạnh rồi đó, ne\dots'' nàng Vu nữ đổ mồ hôi hột, khi cô nàng cảm thấy có những
hơi lạnh đang phả vào người cô. Đã có dấp dáng của những ngôi mộ đầu tiên.

Cậu Tổng cười mỉm, tiếp tục hối thúc mọi người: ``Đây vẫn là buổi chiều mà. Thôi nào,
đi tiếp.''

Cả nhóm bước qua những lối đi ngoắt ngoéo, càng đi sâu càng thấy nhiều bia mộ hiện lên
rải rác giữa những bụi cây. Bầu không khí dần chuyển sang tĩnh lặng, chỉ còn tiếng bước
chân đều đặn và tiếng lá khô vỡ vụn dưới chân.

``Ở đây có vẻ yên bình nhưng\dots\ em cứ thấy hơi lạnh sống lưng sao ấy,'' Shion lên tiếng,
cũng bắt đầu cảm thấy lạnh sống lưng ở nghĩa địa này.

``Đó là do em tưởng tượng ra thôi, Shion-chan\~{}'' nàng Cáo trắng lên tiếng, giọng thốt lên
vẻ tinh nghịch. Lập tức, nàng Phù thủy phản ứng lại: ``Em không phải là trẻ con nha!''

Họ tiếp tục đi, cho đến khi chiếc xe máy của Ryujin hiện ra bên cạnh một ngách nhỏ. Lối
đi vào hẹp hơn, cây cối mọc dày hơn, và những cành gai thấp khiến mọi người phải cúi
người tránh.

``Ái da\dots\ Cây sắc quá\dots'' nàng Vu nữ vừa đi, vừa đụng phải những cây lá gai dọc đường.

Cáo trắng đi phía sau, vừa đi một cách cẩn trọng, vừa dò đường, khẽ hỏi: ``Đến nơi chưa
vậy, Kuon-senpai?''

``Sắp đến rồi,'' cậu con trai đáp lại.

Cuối cùng, cả nhóm dừng chân ở một khoảng đất nhỏ được xây cất khá hoành tráng, nơi
có bảy tấm bia mộ được xếp theo từng hàng ngay ngắn. Không gian nơi đây dường như
càng thêm tĩnh lặng, chỉ còn lại tiếng lá xào xạc và thỉnh thoảng là tiếng chim hót vẳng
xa.

Đây là khu mộ nhà Hikawa, bao gồm hai kỵ nội, cụ ông, cụ bà nội, ông bà nội của bố
Kuon, và một ông là em ruột của ông nội.

``Đây là\dots\ mộ gia đình của anh sao, Kuon-senpai?'' Shion ló mặt nhìn những tấm bia mộ
đó, được khắc tên nắn nót trên đó.

``Ừ. Từ đời cụ tổ của anh đến giờ. Theo anh nhó là đâu đó khoảng\dots\ năm đời,'' cậu nhẩm
lại đời tổ của nhà cậu.

``Tận năm đời sao?'' Fubuki khẽ ngạc nhiên, nhưng không dám to tiếng.

Bố cậu từ bên ngoài bước ra, khẽ gật đầu: ``Đúng như con trai chú nói. Đây chính là khu
mộ của gia đình chú. Gia đình Hikawa thực chất bắt nguồn từ xa xưa nữa cơ, nhưng mà
ở đây bọn chú chỉ nhớ đến năm đời thôi.''

Subaru quay sang cậu Tổng, hỏi với sự tò mò: ``Anh có biết hết tên mọi người không?''

``Không hẳn. Anh chỉ được nghe kể từ các bác trong nhà thôi, vì chưa từng gặp ai trong
số họ cả,'' cậu khẽ lắc đầu, trầm buồn nói. ``Anh cũng không mấy khi ra đây mà.''

Cậu hướng mắt nhìn về các bia mộ với vẻ kính cẩn. Cậu bước tới, lấy trong ba lô ra vài
bó nhang cùng một bó hoa cúc trắng, chuẩn bị cho nghi thức tưởng niệm tổ tiên.

``Được rồi, mọi người. Chúng ta bắt đầu nào.''

Tất cả cùng gật đầu, lần lượt tiến đến đặt những nén nhang lên từng ngôi mộ. Bầu không
khí yên bình nơi đây như thêm phần trang nghiêm, mọi người im lặng cúi đầu, bày tỏ lòng
thành kính với những người đi trước.

Bầu không khí lặng lẽ bao trùm. Gió nhẹ thổi qua, mang theo mùi cỏ cây thoang thoảng,
mấy bông hoa cúc trắng rung rinh trên mộ.

``Mà\dots\ anh nghĩ chúng ta cũng nên làm gì đó cho bọn họ nhỉ?'' Kuon bất giác lên tiếng,
nhìn quanh những tấm bia mộ của tổ tiên mình.

Nói đoạn, Ryujin lấy từ xe máy ra mấy cái chổi cùng bao tay, rồi đưa cho mọi người.
``Mấy đứa phụ chú dọn dẹp chỗ này đi.''

Fubuki cầm lấy chiếc chổi, ánh mắt của cô sáng lên: ``Được thôi, để em biến chỗ này
thành sáng bóng luôn!''

Shion thì tỏ vẻ khó chịu khi cô cũng phải làm công việc vốn dĩ không phải việc của mình,
nhưng vì lòng thành của gia đình mà cô cũng phải miễn cưỡng làm: ``Tự dưng đi dọn
nghĩa trang\dots\ Nhưng thôi, em cũng làm đây.''

Mỗi người một việc. Subaru và Fubuki lo quét lá khô xung quanh, Shion nhặt cỏ dại mọc
lấn ra gần các bia mộ, còn Kuon thay hoa cũ bằng hoa cúc mới mang theo.

Miko nhìn mọi người làm việc, cái cảm giác của nàng vu nữ bên trong cô cũng hối thúc
cô giúp mọi người: ``Em\dots\ cũng muốn phụ mọi người một tay, ne.''

Cậu thấy vậy, đưa cho nàng Vu nữ một bó nhang mới: ``Thế em lấy cái này thắp

Nhận lấy bó nhang, cô nàng khẽ gật đầu: ``Được rồi. Miko sẽ làm tốt nhất có thể!'' Cô
nhẹ nhàng đi tới từng ngôi mộ trong khuôn viên, thắp từng cây nhang đặt trong hũ, rồi
chắp tay cầu nguyện bọn họ.

Sau một lúc, khu vực xung quanh mộ trở nên sạch sẽ hơn hẳn. Không ai lười biếng hay
làm việc gì bất kính, thậm chí cả Shion, vốn hay than thở, hay Fubuki, vốn dĩ rất quậy,
cũng làm tròn nhiệm vụ của mình.

Subaru thở phào nhẹ nhõm sau khi đã xử lý xong đống lá trên khu mộ: ``Xong rồi!'' Cáo
trắng nhìn mọi người với vẻ tự hào: ``Công nhận chúng ta phối hợp tốt ghê.''

Họ đứng lặng yên trước các bia mộ, bố con Kuon lẩm nhẩm vài lời cầu chúc, giọng trầm
thấp nhưng đầy kính trọng.

``Hai người đang nói gì thế, Kuon-senpai?'' nàng Phù thủy trẻ bất chợt hỏi. Kuon sau khi
đã cầu nguyện xong, mỉm cười trả lời: ``Chỉ là cầu mong gia đình được bình an, tổ tiên
phù hộ thôi mà.''

``Miko cũng đã làm điều đó với mọi người rồi mà, ne,'' nàng Vu nữ lên tiếng, không giấu
nổi sự mãn nguyện khi cô vừa làm một điều có ích.

Sau khi dọn xong, họ rời khỏi khu mộ tổ và tiếp tục đi thăm vài ngôi mộ khác gần đó, tiếp
tục cuộc hành trình dọn dẹp mộ mả của gia đình Hikawa.

\dots

\dots

Thế nhưng, ngay khi mọi người vừa rời khỏi, Shion bỗng dưng khựng lại, quay về khu
mộ của tổ tông gia đình Kuon.

Lý do cô ở lại không phải vì chậm chân, mà vì trong lòng có một cảm giác bất an khó tả.
``\textit{Kỳ lạ thật\dots\ Tại sao mình cứ thấy lạnh sống lưng như vậy chứ? Có phải do tưởng tượng
  không?}''

Cô đứng lặng một lúc trước những bia mộ của gia đình Hikawa, trong lòng cô nổi lên một
sự bất an. Gió bỗng thổi mạnh hơn, làm cây cối xung quanh phát ra những âm thanh xào
xạc đầy u ám.

``\textit{Không biết có chuyện gì xảy ra không nhỉ\dots\ Nhưng cảm giác này\dots\ khó chịu thật đấy.}''

Như bị thúc giục bởi một nỗi sợ mơ hồ, cô quyết định làm điều mà ngay chính bản thân
cũng cảm thấy kỳ quặc. Lấy một cái xẻng nhỏ giấu trong túi thần bí của mình (thứ cô
luôn mang theo nhưng chẳng ai biết lý do), Shion âm thầm đào một cái huyệt ngay trong
khuôn viên mộ gia đình Hikawa.

Cô chọn một khoảng đất vừa đủ, ước chừng chiếc huyệt mà cô định đào bằng những chiếc
quanh đó. Cô đào ra một cái hố vừa đủ lớn, vừa đủ sâu để gây cảm giác ``đề phòng,'' như
thể chuẩn bị sẵn cho một điều chẳng lành mà chính cô cũng không thể giải thích rõ ràng.

``\textit{Chẳng biết mình làm thế này có đúng không\dots\ Nhưng nếu có chuyện xảy ra, ít nhất cũng
  có chỗ để\dots}'' Câu nói dở dang, cô lắc đầu thật mạnh, như muốn xua tan ý nghĩ đen tối
đang nhen nhóm trong đầu.

Khi cái huyệt đã xong, Shion nhanh chóng phủ lớp đất mỏng lên trên để giấu đi dấu vết,
đồng thời nhặt một ít cỏ khô và đá nhỏ để nguỵ trang. Cô đứng dậy, phủi tay, nhìn quanh
một lượt để chắc chắn không ai phát hiện.

``\textit{Không được nói với ai cả\dots\ Nếu để họ biết, nhất là Kuon-senpai, chắc chắn sẽ bị coi là
  điềm xấu}.''

Cô quay người rời đi, bước chân nhanh hơn bình thường như để tránh ánh mắt vô hình
nào đó đang dõi theo. Trước khi tiến tới khu rừng cây rậm rạp, Shion ngoái lại một lần
cuối, nhìn về phía những bia mộ và cái huyệt mà cô vừa đào.

Cô thì thầm, gần như cầu nguyện: ``Mong rằng phỏng đoán mình là sai\dots''

Cảm giác bất an vẫn không rời khỏi Shion, nhưng cô quyết định giữ tất cả những gì đã
làm trong bí mật, không hé lộ với bất kỳ ai, dù chỉ một lời.

\ct

Trên con đường đi thăm những ngôi mộ khác vào buổi chiều mùng Ba Tết, họ vừa được
Ryujin nghe kể những câu chuyện kỷ niệm về họ, đồng thời cũng được Kuon giải thích
về vai vế trong gia đình của nhà Hikawa.

Càng lúc, họ càng được hiểu sâu hơn về gia thế của đại gia đình cậu Tổng này.

Những ngôi mộ mà họ đến thăm cũng thuộc dạng ``độc nhất vô nhị,'' khi một trong số đó
thật sự độc đáo: mộ của cụ nội bên ngoại của Kuon lại được đặt ngay giữa\dots\ một cánh
đồng lúa. Cảnh tượng khiến cả nhóm đứng lại nhìn, vừa bất ngờ vừa buồn cười.

Cậu Tổng lên giọng châm chọc: ``Có khi nếu em chết thì bọn anh sẽ đặt mộ em ở giữa hồ
đấy, Subaru-chan.''

Nàng Vịt lập tức phản ứng, khuôn mặt đỏ bừng, phồng má lên: ``Mồ! Đừng có trù ẻo em
thế chứ! Với cả em không phải là vịt!''

Cáo trắng bật cười khúc khích, còn nàng Vu nữ hùa theo: ``Em đang là vịt còn gì. Đặt mộ
giữa hồ thì hợp với em lắm luôn ấy, ne!''

Subaru hét lên, mặt đầy vẻ tức tối vì bị khịa: ``\textbf{Lại còn cả mấy người nữa! Quá đáng!}''

Shion nhướng mày, đứng khoanh tay quan sát, rồi mỉm cười đầy ẩn ý: ``Biết đâu đây là
điềm báo gì thì sao? Subaru, bà nên cẩn thận hơn đấy.''

Nàng Vịt quay qua nhìn cô với ánh mắt bất lực: ``Bà thì cứ thêm dầu vào lửa thôi! Người
gì đâu độc mồm độc miệng ghê luôn!''

Nàng Phù thủy nhún vai, đôi mắt sắc sảo liếc qua Kuon: ``Thật ra, đặt mộ giữa đồng thế
này cũng có phong thủy tốt đấy. Nhưng mà\dots\ ở giữa hồ thì chắc chắn sẽ đặc biệt hơn.''

Cậu Tổng bật cười: ``Anh nói đùa thôi mà sao mọi người làm căng thế! Được rồi, thôi tập
trung dọn dẹp nào.''

Họ cũng làm tương tự với mộ của cụ ấy, mỗi người một tay một việc. Những bụi cỏ dại
bị nhổ sạch, vài nhành hoa tươi được cắm vào bình, và không khí trang nghiêm trở lại.

Khi mọi việc đã xong, cả nhóm tất tả ra về. Riêng bố Kuon vội vã quay về xưởng để ``mở
hàng đầu năm'', còn Kuon và các cô gái cùng trở lại nhà để chuẩn bị cho những hoạt động
cuối cùng của ngày lễ.

\ct

Tại phòng thờ gia đình Hikawa, ngay khi đi thăm mộ về, Kuon sắp mâm ngũ quả lên bàn
thờ rồi thực hiện thờ cúng, giống như gia đình cậu hay làm vào mọi năm.

Miko đứng phía sau, chăm chú quan sát từng hành động của cậu. Những món trái cây đủ
màu sắc được xếp ngay ngắn trên khay, mỗi loại mang một ý nghĩa riêng. Thấy thế, nàng
Vu nữ nhẹ nhàng hỏi: ``Cái đó có phải là thờ phụng tổ tiên không, ne?''

Kuon vừa cúng xong, quay lại gật đầu, đáp: ``Đúng rồi đấy. Anh nghĩ Miko-chan cũng
làm thế hả?''

Miko chắp tay, ánh mắt hơi bối rối, nhưng cũng đầy hoài niệm: ``Hồi trước, ở Ngôi đền
ảo, em cũng làm những nghi thức kiểu này. Nhưng\dots\ có lẽ không đầy đủ và trang trọng
được như thế này.''

\ct

Cậu mỉm cười, tỏ vẻ hơi ngạc nhiên: ``Anh cũng không biết em từng làm ở ngôi đền đấy.''
Cậu đẩy khay mâm ngũ quả lên ngay ngắn hơn, rồi tiếp tục: ``Đền thờ chắc cũng có cách
làm riêng, phải không? Anh nghĩ việc mình thành tâm là điều quan trọng nhất, dù cách
làm có hơi khác nhau.''

Miko gật đầu, đôi mắt ánh lên sự cảm kích. Cô bước tới gần, chỉ tay vào từng loại quả
trên mâm, tò mò hỏi: ``Nhưng mà\dots\ tại sao lại là năm loại quả này, ne? Có phải có ý
nghĩa gì đặc biệt không?''

Kuon cười, cùng với cô ngồi xuống và kiên nhẫn giải thích: ``Cũng tùy từng vùng miền
thôi em ạ. Chung quy lại, chúng tượng trưng cho mong ước về sự đầy đủ, hạnh phúc, phú
quý, và những điều tốt đẹp khác.''

``Ví dụ, ở đất Bắc bọn anh, mâm ngũ quả sẽ có chuối xanh, bưởi, cam, và nho, tượng trưng
cho ngũ hành,'' cậu nói, bắt đầu chỉ vào từng quả trên mâm.

Nải chuối xanh ứng với Mộc (màu xanh), tượng trưng cho sự sum vầy, đầm ấm, đồn thời
nó còn có ý nghia che chở khi nó đóng vai trò làm giá đỡ trên mâm ngũ quả.

Bưởi trắng ứng với Kim (màu trắng), biểu trưng cho sự viên mãn, vẹn tròn, gửi gắm mong
ước về một năm mới hanh thông.

Phật thủ ứng với Thổ (màu vàng), tượng trưng cho sự may mắn, tài lộc.

Cam ứng với Hỏa (màu đỏ), tượng trưng cho sự thành đạt, thăng tiến trong sự nghiệp.

Nho ứng với Thủy (màu đen), biểu tượng của sự sinh sôi và nảy nở, ngoài ra còn đại diện
cho sự thành công.

Sau khi chăm chú nghe những gì cậu giải thích, nàng Vu nữ gật gù, bật cười nhẹ: ``Nghe ý
nghĩa thật đấy. Anh nói rằng mỗi vùng có một kiểu sắp xếp khác nhau, vậy nó khác nhau
như thế nào vậy, ne?''

``Đúng thế. Ở miền Nam thì thường tránh chuối, vì chữ `chuối' nghe giống `chúi', mang
hàm ý không may mắn.''

Miko đưa tay lên miệng, che nụ cười khúc khích: ``Thú vị thật đấy ne. Em không biết ở
đây lại nhiều khác biệt như vậy.''

Cậu Tổng đứng dậy, đặt tay lên khay mâm, nhìn nàng Vu nữ đầy chân thành: ``Anh nghĩ
em có thể mang những nét văn hóa này giới thiệu với fan của em, nếu họ chưa biết. Chắc
chắn họ sẽ rất thích.''

Miko ngẫm nghĩ một lúc, rồi đáp lại bằng nụ cười tươi rói: ``Em sẽ thử xem! Nhưng mà\dots
liệu em có làm tốt được như anh không nhỉ?''

Cậu bật cười: ``Đến PON như em gái anh còn có thể làm được mà. Chỉ cần làm bằng tấm
lòng là được rồi, Miko-PON.''

Nàng Vu nữ phồng má lên phụng phịu: ``Đừng có so em với cô ấy nữa, ne. Em không
thích bị gọi như vậy đâu.''

Không khí trong phòng thờ trở nên nhẹ nhàng, ấm cúng. Những lời chia sẻ và câu chuyện
của họ như tô thêm sắc màu cho ngày Tết, nối dài truyền thống từ thế hệ này sang thế hệ
khác.

Sau khi họ bày mâm ngũ quả và thực hiện nghi lễ xong, hai người họ ngồi xuống, tiếp tục
trò chuyện với nhau về kỷ niềm mà Sakura Miko đã từng làm trong Ngôi đền ảo năm nào.

Căn phòng thờ yên tĩnh, chỉ còn lại ánh sáng từ ngọn nến và hương trầm tỏa nhẹ trong
không khí. Kuon bắt đầu hỏi, với sự hiếu kỳ: ``Mà, em nói rằng em từng làm ở Ngôi đền
ảo đúng chứ? Em thường làm gì ỏ đó?''

Miko chống cằm, ánh mắt hơi xa xăm, đáp: ``Miko thường lau dọn đền thờ thôi. Ít nhất
là trước khi em gia nhập \hll, ne.''

``Nghe giống kiểu công việc vừa vất vả vừa yên bình nhỉ. Nhưng chắc cũng không thiếu
những kỷ niệm đặc biệt đâu, đúng không?'' cậu cười nhẹ.

Nàng Vu nữ gật đầu, rồi bất chợt mỉm cười: ``Đúng thế! Em từng sống với Kintoki và
Thần ở đó.''

Nghe những tên lạ hoắc, cậu Tổng nhướng mày tò mò: ``Kintoki? Thần? Họ là ai vậy?''

Nàng Vu nữ chậm rãi giải thích, đôi mắt long lanh lên khi nhắc về những ký ức cũ:
``Kintoki là con mèo hồng biết nói của em. Còn ông Thần ấy thì\dots\ thôi, bỏ đi.'' Cô nàng
khẽ bật cười khi nhớ về những người đã từng gắn bó với cô suốt khoảng thời gian đó.

Kuon nhíu mày, nhưng ánh mắt lại pha chút hài hước: ``Ồ, anh cá chắc là họ nghiêm túc
lắm luôn.''

Miko khẽ cười, nhưng tiếng cười ấy nhanh chóng trở nên trầm lắng hơn khi cô thầm nhớ
lại những kỷ niệm xưa cũ: ``\dots Đúng vậy. ne. Em với Kintoki đã biết nhau từ khi em còn
bé mà ne. Nó thậm chí còn coi em như em gái.''

Cậu Tổng cười lớn, châm chọc: ``Vậy là em đã có người chăm sóc từ trước rồi. Anh cũng
chỉ về nhì thôi.'' Nàng Vu nữ giật mình, mặt đỏ bừng, vội vàng xua tay: ``Ý-Ý em không
phải thế! Mà tại sao anh nói thế chứ!?''

``Nói đùa thôi mà. Nhưng cuộc sống trong đền chắc cũng tẻ nhạt nhỉ?'' cậu cười xòa, rồi
quay về nói với sự thoải mái.

Nàng Vu nữ ngừng lại một chút, rồi đáp: ``Ngoài Kintoki và Thần ra, em không có cơ hội
nói chuyện với ai cả. Nhưng em nhớ\dots\ Roboco-san đã từng đến chơi một lần rồi\footnote{Xem {\flaregothic ÉPISODE -4}, quyển Tiền truyện.}.''

Kuon nghiêng đầu, thắc mắc: ``Roboco-chan sang chơi à? Em ấy làm gì ở đó?'' Miko
mỉm cười, ánh mắt phảng phất nét ngây thơ khi nhớ lại: ``Em với cậu ấy chỉ chơi mấy trò
ngớ ngẩn thôi. Hồi đó em còn chưa thể tự tin bắt chuyện với cậu ấy như bây giờ được.''

``Hồi đó chắc em ngại giao tiếp lắm nhỉ.''

``Vâng, hồi đó là như thế ạ, ne. Nhưng bây giờ\dots\ em thấy mình đã thay đổi rất nhiều rồi.''

Nàng Vu nữ bắt đầu kể lại những ký ức mà cô nàng đã từng sống và sinh hoạt ở Ngôi đền
ảo, nơi giờ đây chỉ còn lại trong tâm trí. Cô nhắc đến sự biến mất của Kintoki và Thần
khi cô rời đền, gia nhập \hll\ và bước vào một thế giới mới.

Có những kỷ niệm mà khi cô nàng kể thì khiến cho ai nghe cũng cười rơi nước mắt, nhưng
cũng có những câu chuyện khi nghe xong cũng khiến ta đồng cảm.
Chẳng hạn, khi cô nàng đã phải tự luyện tập giao tiếp với mọi người\footnote{Xem {\flaregothic ÉPISODE -5}, quyển Tiền truyện.}, đánh vật với cây
thông hậu Giáng Sinh\footnote{Xem {\flaregothic ÉPISODE -7}, quyển Tiền truyện.}, hay cùng Kintoki luyện tập tặng sô-cô-la ngày Lễ tình nhân\footnote{Xem {\flaregothic ÉPISODE -1}, quyển Tiền truyện.}\dots\ Đó là những kỷ niệm mà có lẽ đã khắc sâu vào trong lòng của nàng Vu nữ này.

Cô nói thêm, giọng chùng xuống: ``Bản thân Kintoki mà em có bây giờ chỉ là một bản tái
tạo lại của bản gốc thôi. Nó không còn nghiêm túc và chăm chỉ như trước được nữa.''

``Và nó cũng không nói được nữa à?'' cậu nhìn cô nàng đang đượm buồn ấy, tỏ vẻ đồng
cảm với cô. Cậu nhẹ nhàng đặt vai lên Miko-chi, an ủi: ``Nhưng điều quan trọng là những
ký ức ấy vẫn luôn sống mãi trong em. Và dù thế nào đi nữa, anh tin rằng Kintoki, hay
cả ông Thần ấy, đều tự hào về em. Họ sẽ rất vui khi được nhìn thấy một cô gái bé bỏng
trưởng thành đấy, Miko-chan ạ.''

Miko ngẩng đầu lên, ánh mắt ngập tràn sự cảm kích. ``Cảm ơn anh, Kuon-senpai. Đã lâu
rồi em mới có cơ hội giãi bày tâm sự như vậy đó, ne.''
Cô nàng Vu nữ bật dậy, giương ra bộ mặt tươi tắn vốn có của một thần tượng: ``Thôi nào,
để em phụ anh dọn dẹp lại bàn thờ và sắp xếp phòng ốc ne.''

Kuon bật cười, bước theo cô: ``Được thôi. Anh sẽ không từ chối đâu, Miko-chi.''

Căn phòng thờ một lần nữa trở nên bận rộn, nhưng không khí lại đầy ắp tiếng cười và sự
ấm áp. Những câu chuyện xưa cũ và hiện tại như hòa quyện, tạo nên một bức tranh đầy
cảm xúc trong những ngày đầu năm mới.

\ct

\begin{center}
  \uline{3:00 chiều\\Đền \ruby{Fuwan}{Hùng Vương}\ - \ruby{Fushuu}{Phú Thọ}}.
\end{center}

Kuon cùng các thành viên của \hll\ đến dự lễ hội Tết tại một ngôi đền cổ kính, ở một vùng đất được coi là tổ tiên của người Etsunan - đền Hùng Vương, hay đền Vương. Đây là một trong những đền thiêng nhất của người dân tại xứ này.

Matsuri xoay người một vòng, đôi mắt sáng rỡ nhìn khắp nơi: ``Nơi này rộng quá đi!
Đúng kiểu lễ hội mà Matsuri em đây yêu thích!''

Bầu không khí lễ hội đang trở nên náo nhiệt hơn bao giờ hết. Những dãy
đèn lồng đỏ rực trải dài từ cổng đền vào sâu bên trong, lung linh dưới ánh mặt trời. Tiếng
trống, tiếng chiêng hòa lẫn với âm thanh của những màn múa lân, tạo nên một bầu không
khí buổi chiều ngày Mùng Ba vừa trang nghiêm, vừa sôi động.

Choco nhanh chóng kéo tay nàng Cổ độnh lại, lo lắng nhắc nhở: ``Bình tĩnh nào, Matsurisama.
Đừng có chạy lung tung kẻo lạc đấy.''

Ayame bước tới gần, mắt chăm chú nhìn những đèn lồng đang lập lòe sáng: ``Hơi giống
lễ hội mà em từng tham gia trước đây nha-dazo. Nhưng ở đây đông hơn nhiều.''

Kuon khoanh tay, nhìn các cô gái đang tản ra khám phá ở ngôi đền này, cười nhẹ: ``Có lẽ
các em chưa biết, nhưng Fushuu là một trong những nơi tổ chức lễ hội Tết rầm rộ và lâu
đời nhất của cả vùng Etsunan đấy. Tất nhiên chưa là gì so với ngày Giỗ Tổ đâu.''

Mel gật đầu đồng tình, thêm vào: ``Nhưng mà, em nghe nói các thành phố khác quanh
đây tổ chức còn hoành tráng hơn. Ở đây đã đẹp rồi, chắc mấy chỗ đó còn tuyệt hơn nữa.''

Nàng Lễ hội nghe nàng Dơi giải thích vậy, lập tức phản ứng, giọng đầy phấn khích: ``Vậy
tại sao chúng ta không đi mấy chỗ đó luôn!?''

Cậu Tổng bật cười, xua tay đáp: ``Thứ nhất, nó xa lắm. Thứ hai, các em đi một lần ở đây
là đủ hình dung được rồi. Mỗi nơi có nét đặc sắc riêng, nhưng lễ hội ở nơi này cũng sôi
động chẳng thua kém đâu.''

Kuon gật đầu: ``Đúng vậy. Đây là nơi tinh hoa văn hóa ngày Tết hội tụ. Người ta đến đây
không chỉ để vui chơi mà còn để giữ gìn những phong tục cổ truyền nữa.''

Aqua ngẩng đầu, hỏi ngay: ``Những hoạt động trải nghiệm đó là gì vậy, Kuon-senpai?''

Kuon khẽ mỉm cười khi Aqua hỏi, đôi mắt ánh lên sự thích thú. Không khí lễ hội xung
quanh thật rộn ràng với tiếng trống rộn rã, tiếng chiêng vang vọng và mùi hương nhang
trầm phảng phất. Trẻ con chạy nhảy tung tăng, người lớn thì mặc áo dài truyền thống, trò
chuyện vui vẻ trong khi mọi người cùng nhau trải nghiệm các hoạt động văn hóa.

``Trước tiên thì, anh muốn mời những người thuộc loại ba từ bài kiểm tra toán ra đây,''
cậu nói.
Nàng Hầu gái nhăn mặt, lập tức bước lên: ``Anh định giở bài gì với bọn em vậy?'' Subaru
cũng chen vào, tay chống hông: ``Đừng có bắt bọn em làm cái gì quá sức đấy nha!''

Kuon bật cười, chỉ tay lần lượt vào Aqua, Ayame, Subaru, Miko và Okayu: ``Không cần
phải lo đâu. Anh đã chuẩn bị sẵn hoạt động vừa sức cho mấy đứa rồi.''

Ayame nhìn cậu với vẻ mặt cảnh giác, nhưng rốt cuộc vẫn bước tới. Okayu thì chỉ khẽ
thở dài, đôi mắt như đang ngầm nói rằng cô đã đoán được phần nào ý định của cậu Tổng.
Dẫu biết rằng những hoạt động ngày Tết tại đền chùa luôn mang tính truyền thống và nghi
lễ, họ vẫn không giấu nổi sự hồi hộp.

Cả nhóm tập trung xung quanh Kuon, trong khi những người khác đứng gần đó theo dõi,
vẻ mặt vừa tò mò vừa buồn cười. Kuon, tuy vẫn giữ thái độ đùa cợt, khẽ nghiêng người
xuống thì thầm điều gì đó khiến nhóm người kia đồng loạt nhìn nhau, rồi quay lại nhìn
cậu với ánh mắt ngờ vực\dots

\ct

Tất cả mọi người xếp thành hai hàng, hòa vào dòng người tiến về phía tượng Phật lớn
trong khuôn viên đền. Từng người, từng người theo đoàn bước đi trên những bậc cầu
thang dài tưởng như vô tận. Không khí trang nghiêm dần bao trùm, nhưng vẫn không
thiếu những tiếng thì thầm nho nhỏ xen lẫn tiếng cười khúc khích.

Mọi người, ngoại trừ A-chan và Kuon, đều mặc áo dài - bộ trang phục truyền thống của
Etsunan - để phù hợp với không khí ngày Tết và thể hiện sự tôn trọng khi vào đền.

Shion kéo vạt áo dài của mình, gương mặt không giấu nổi sự khó chịu: ``Tại sao em phải
mặc bộ quần áo này chứ\dots\ Ngứa quá đi!''

Roboco, đứng cạnh, cũng tỏ vẻ bất bình không kém: ``Trông hơi kỳ quặc kiểu gì ấy. Không
quen chút nào.''

Kuon, quay lại gõ nhẹ vào trán hai cô gái đang than vãn ấy: ``Yên nào. Đây là nơi linh
thiêng, giữ phép lịch sự tối thiểu chứ.''

Mel nhướng mày, khoanh tay trước ngực: ``Chị thấy hai người thường ngày mặc hở hang
nhất trong nhóm mà.''

A-chan lập tức chen vào, ánh mắt như muốn bắt lỗi cô ấy: ``Cả cô cũng thế đấy, Yozora-san.
Bộ quần áo mặc thường ngày của mấy người mới thực sự kỳ lạ trong hoàn cảnh này.
Vào đền mà mặc đồ không liên quan tới Tết là thiếu nghiêm túc đấy.''

Miko ngơ ngác, vội vàng hỏi: ``Ể? Nhưng tôi tưởng là ta có thể mặc tự do được mà?''
Nàng Đeo kính thở dài, ôn tồn giải thích: ``Về lý thuyết là thế. Nhưng tôi muốn mọi người
nghiêm túc hơn một chút trong dịp này. Tôi và Hikawa-san đã mất công thuê trang phục
truyền thống cho tất cả rồi mà.''

Kuon khẽ gật đầu, chốt lại bằng giọng điệu kiên quyết: ``Ở đây không giống như ở Nhật
đâu. Nhập gia thì phải tùy tục chứ, đúng không?''

Cả nhóm im lặng trong giây lát, rồi tiếp tục bước về phía trước, hòa vào dòng người, tiếp
tục cuộc hành hương tới bức tường Phật.

Ở đền Vương, các hoạt động lễ hội thường gắn liền với ý nghĩa ``hành hương về cội
nguồn,'' nơi mọi người cùng nhau thắp hương tưởng nhớ các Vua Hùng và cầu mong may
mắn, bình an cho năm mới. Cảnh người dân tấp nập nối dài thành từng đoàn, từ từ tiến về
phía khu vực chính điện để dâng lễ, tạo nên không khí trang nghiêm nhưng không kém
phần ấm áp.

Aqua và Subaru đứng ở đầu đoàn, với tay của nàng Vịt đang cầm bát hương. Nàng Hầu
gái than phiền, tay phẩy phẩy: ``\hl{Khói bay mù mịt quá\dots}''

Subaru cũng tỏ vẻ khó chịu không kém, nhưng lai nhẫn nhịn hơn: ``Cũng hơi khó chịu
nữa. Cơ mà, nó cũng là trải nghiệm đáng thử đấy.''

``\hl{Tôi có thể cầm được không?}'' cô nàng tỏ vẻ hiếu kỳ, với bát hương đang nghi ngút khói.
Lập tức, nàng Vịt nhăn mặt, khẽ lắc đầu: ``\hl{Thôi bà ơi. Đến cả lau dọn bà còn rơi vỡ loảng
  xoảng, làm sao mà tôi tin tưởng đưa cho bà được? Huống chi đây lại bát hương nữ\dots}''

Cảm thấy như đang bị xúc phạm vào tinh thần hầu gái của mình, Aqua thốt lên: ``\hl{Ể\\~{}{}?
  Không công bằng!}''

Kuon quan sát đoàn người đang lặng lẽ tiến bước, rồi nhìn về phía Aqua và Subaru vẫn
đang lục đục với bát hương: ``\hl{Hai đứa trật tự đi cho anh nhờ cái\dots}''

Aqua quay sang, khẽ phồng má: ``\hl{Nhưng tại sao chúng ta phải đi chậm thế này? Không
  phải cứ dâng hương xong là xong sao?}''

Nàng Vịt lườm người đồng hành một cái, tỏ vẻ không hài lòng với phản ứng đó: ``\hl{Bà
  không biết đây là hành hương à? Đây là một phần nghi thức truyền thống, bà phải nghiêm
  túc vào}!''

Kuon mỉm cười, giải thích thêm: ``\hl{Đúng vậy. Hành hương về đây không chỉ là để dâng
  hương mà còn là để mọi người suy ngẫm về cội nguồn, về những giá trị mà các thế hệ
  trước đã truyền lại. Cả hai đứa cũng nên trân trọng điều này hơn.}''

Nàng Hầu gái im lặng một chút, ngẫm nghĩ lại một lúc lâu về ý nghĩa của cuộc hành
hương. Cô khẽ gật gù, hiểu được những gì mà cậu vừa giải thích: ``\hl{Ừm\dots{} nghe vậy thì
  cũng thấy ý nghĩa thật đấy. Nhưng Subaru vẫn không được cầm bát hương lâu hơn em
  đâu!}''

Subaru ngán ngẩm, nhưng rốt cuộc vẫn khẽ bật cười: ``\hl{Được rồi, vậy bà cầm đi. Nhưng
  nhớ đừng làm rơi đấy, Rơi-vỡ-chan!}''
Nàng Hầu gái phụng phịu phồng má, cầm chiếc bát hương từ tay nàng Vịt.

\ct

Cuộc hành hương dẫn mọi người đến một ngôi đền nhỏ nằm sâu bên trong khuôn viên, nơi
đặt bức tượng Phật uy nghiêm được trang trí bằng những vòng hoa rực rỡ. Dòng người
lần lượt tiến vào, ánh mắt ai cũng ánh lên vẻ thành kính khi đặt lễ vật và thắp hương trước
tượng.

Aqua và Subaru, với bát hương trên tay, cũng bước lên trước bàn thờ để dâng lễ. Sau khi
cúi chào và cẩn thận đặt bát hương xuống, cả hai cùng lui về phía sau, nhập vào đoàn
người đang lặng lẽ cầu nguyện.

Nàng Hầu gái khẽ thì thầm, ánh mắt lo lắng lướt qua hàng tá người đang đứng cùng mình:
``\hl{Nhiều người đến dự quá\dots}''

Nàng Vịt đáp lại, giọng tỏ vẻ bình thản như ``cân đường hộp sữa'': ``\hl{Đương nhiên. Đây là
  lễ hội mà.}''

Aqua nhìn dòng người đông đúc, mắt xoáy như chong chóng: ``\hl{L-Làm sao bây giờ\dots{} Căng thẳng quá đi mất\dots}''

``\hl{Tưởng tượng bà đang ở trên sân khấu ấy,}'' Subaru cố gắng trấn an bạn mình. Nhưng có vẻ như ``Hành tím'' vẫn đang rất lo lắng: ``\hl{Tôi cảm thấy nơi này khác hoàn toàn so với ở trên sân khấu á\dots}''

``\hl{Vậy cứ coi họ là rau củ là được.}''

Nghe lời Subaru, Aqua nhắm mắt, cố gắng hình dung: ``\textit{Khán giả là củ cải\dots\ bắp cải\dots\ cà
  rốt\dots}'' Khi cô mở mắt lại, trước mắt cô vẫn là dòng người kia đang im lặng cầu nguyện.

Có vẻ như chiêu bài mà ``Ho'' và ``Chim nhỏ'' bày ra cho ``Biển cả''\footnote{Liên hệ tới một cảnh trong bộ anime \textit{Love Live! School idol project}.} không hiệu quả lắm
với cô nàng. Cô nhắm tịt mắt lần nữa, cố gắng thử lại một lần nữa.

Nhưng rồi, có lẽ do quá căng thẳng mà trong một phút quá cao hứng, nàng ``Hành tím''
bất ngờ bật dậy, hét lớn: ``\textbf{\color{red}Khán giả là rau củ!}''

``Aqua-chan!?'' Subaru thảng thốt khi nghe cô bạn hét lớn như vậy.

Không khí yên tĩnh trong đền phút chốc bị phá tan. Những ánh mắt của mọi người ngạc
nhiên từ xung quanh lập tức đổ dồn về phía cô, khiến Aqua đỏ bừng mặt.

Cả không gian bên trong ngôi đền nhỏ như dừng lại trong một khoảnh khắc im lặng đầy
ngỡ ngàng. Sau đó, những tiếng xì xào bắt đầu râm ran trong đám đông.

``{\color{red} Cô gái đó vừa nói gì thế nhỉ?}'' ``{\color{red} Khán giả là rau củ? Ý cô ấy là sao vậy?}'' ``{\color{red} Mấy người
  trẻ tuổi bây giờ lạ thật\dots\ Nhưng trông cô nàng ấy đúng là dễ thương.}''

Một người phụ nữ trung niên đứng gần đó cười nhẹ, nói với chồng mình: ``{\color{red} Chắc cô ấy
  căng thẳng quá nên nói linh tinh thôi.}''

Ở một góc khác, một ông lão với chiếc nón lá khẽ vuốt râu, nửa trêu chọc, nửa tò mò:
``{\color{red} Đến đền mà nghĩ đến rau củ\dots\ Đúng là đầu óc trẻ con bây giờ sáng tạo thật.}''

Nghe những lời xì xào xung quanh, nàng Hầu gái hậu đậu chỉ biết cúi gằm mặt, cố nén sự
xấu hổ: ``Trời ơi\dots\ Tại sao mình lại hét lớn như vậy chứ\dots''

Subaru đứng bên cạnh thì không thể nhịn nổi cười, dù tỏ vẻ lo lắng cho cô bạn. Cô ghé
xuống tai của bạn mình và thì thầm châm chọc: ``Bà nổi tiếng rồi đấy. Người ta sắp gọi
bà là `thánh rau củ' mất thôi!''

Aqua ngẩng đầu, bối rối lắc đầu, nhỏ giọng phản bác lại: ``Không phải vậy! Không phải
vậy đâu! Đừng nói nữa mà, Subaru-chan\dots''

Pekora, đang chăm chú dâng hương, từ phía dưới ngước lên, búng tay ra hiệu về phía hai
người: ``Này! Aqua-senpai, không cần lo đâu, peko! Có ai giận chị đâu. Nhưng nhớ là
giữ trật tự đi, peko!''

Kuon và A-chan, đứng gần đó, chỉ biết lắc đầu bất lực. Nàng Đeo kính thở dài, thấp giọng
nhắc nhở: ``Hikawa-san, lần sau cậu chọn người cầm hương cẩn thận hơn nhé.''

Kuon nửa đùa nửa thật, mỉm cười: ``Đúng là một buổi lễ hội sống động hơn mong đợi
nhỉ?''

Cậu tiến tới dòng người phía trước, luồn qua dòng người đang bàn tán xôn xao kia, tiến
tới cô nàng Hầu gái đang ``xì khói như bình siêu tốc'' vì xấu hổ. Cậu lặng lẽ đặt tay lên vai cô, giọng điệu nhẹ nhàng nhưng không kém phần nghiêm túc: ``Lần sau, nhớ suy nghĩ
trước khi nói nhé. Đây là nơi trang nghiêm, không phải sân khấu đâu.''

Nàng Hầu gái ngẩng lên, đôi mắt rưng rưng, lí nhí đáp: ``Em\dots\ em xin lỗi\dots\ Em không
cố ý đâu\dots''

Mọi người trong nhóm cũng nhìn Aqua với ánh mắt vừa thương cảm, vừa buồn cười.
Ayame nhún vai, lẩm bẩm, nở một nụ cười gian: ``Đúng là Akutan\dots\ Không bao giờ làm
ai thất vọng trong việc gây rắc rối mà.''

Cả nhóm bật cười nhẹ, xoa dịu bầu không khí đang có phần căng thẳng. Dù sự cố nhỏ này
khiến Aqua ngại ngùng, nhưng nó lại vô tình trở thành một khoảnh khắc đáng nhớ khác
cho mọi người trong chuyến hành hương ngày Tết này.

\ct

Sau cuộc hành hương vất vả và đầy ý nghĩa, cả nhóm tiếp tục di chuyển đến một khu vực
khác trong lễ hội. Tiếng trống dồn dập vang lên hòa quyện với tiếng hò reo từ đám đông,
báo hiệu một màn biểu diễn đặc sắc.

Họ dừng chân trước một sân khấu nhỏ được dựng tạm, với vài bảy, tám chiếc trụ có độ
cao khác nhau, nơi một nhóm bốn người đang khéo léo điều khiển từng động tác, tạo nên
hình ảnh sống động của một con lân đầy sắc màu. Bộ lông đỏ vàng rực rỡ của con lân
đung đưa nhịp nhàng theo từng tiếng trống, thu hút ánh nhìn của mọi người.

``Hoá ra con lân đó có người điều khiển ở dưới à,'' Haato thốt lên, đôi mắt ánh lên vẻ ngạc
nhiên.

``Chính xác ra là họ đang đóng giả làm lân đấy,'' Kuon nhẹ nhàng giải thích, nụ cười
thoáng trên môi.

``Thế mà tôi cứ tưởng họ đi múa với con lân thật chứ\dots'' Matsuri lẩm bẩm, vẻ thất vọng
hiện rõ trên khuôn mặt. Dù vậy, cô vẫn thở phào nhẹ nhõm thì ít nhất cô không phải nhìn
thấy con lân thật.

Rushia chen ngang, vừa nói vừa tưởng tượng cảnh mình làm bùa để điều khiển con lân
thật: ``Như thế thì em sẽ có nhiều việc để làm rồi-nanodesu.''

``Đúng đó, với lại-'' Choco quay sang, định thêm ý kiến thì chợt khựng lại. Cô liếc mắt
nhìn quanh, gương mặt lộ vẻ hoang mang. ``Ủa. Ayame-sama đâu rồi?''

Cả nhóm im lặng vài giây, rồi như chợt nhận ra điều gì đó, ánh mắt họ đồng loạt hướng
về phía sân khấu. Không nằm ngoài dự đoán, nàng Quỷ sừng nhỏ bé đã lẻn ra phía sau
đám đông và nhanh nhẹn chạy đến chỗ nhóm múa lân. Trên khuôn mặt Ayame là vẻ phấn
khích rõ rệt.

``Em ấy lại làm trò gì nữa đây?'' Kuon thở dài, bước tới gần sân khấu.

Lúc này, Ayame đã đứng đối diện với nhóm múa lân, dõng dạc xin phép với sự hớn hở,
đôi mắt sáng rực: ``{\color{red} Nè nè, Con người-sama, cho tôi thử được không? Tôi muốn thử điều
  khiển con lân này!}''

Đám đông dường như bị cuốn hút bởi sự tự tin bất ngờ của nàng Quỷ sừng. Một thành
viên trong đội tỏ vẻ ngạc nhiên khi có một người ngoài như cô muốn xin phép biểu diễn:
``{\color{red} Ể? Cô có thể làm được không}?''

Nàng Quỷ sừng mỉm cười một cách tự tin: ``{\color{red} Tôi tin là tôi làm được!}'' Nhìn với vẻ mặt
sáng láng của cô nàng, cậu thanh niên đó cuối cùng cũng gật đầu, giao lại phần đầu lân
cho Ayame.

``{\color{red} Cô có chắc là cô có thể thực hiện được chứ? Đây không phải là thứ mà những người
  nghiệp dư có thể làm được đâu\dots}'' một thành viên trong đội múa lân vẫn tỏ vẻ lo ngại.

Nàng Oni giơ ngón tay cái hướng lên trên, mỉm cười: ``{\color{red} Yên tâm đi! Tôi đã từng tham dự
  một đội múa lân như vậy từ trước rồi mà, nên tôi thạo lắm!}''

Tiếng trống lại vang lên, và màn múa lân phiên bản Oni chính thức bắt đầu. Với thân hình
nhỏ nhắn và sự linh hoạt trời ban, Ayame điều khiển phần đầu lân vô cùng sống động.
Mỗi cú vung đầu hay cú nhảy đều được đám đông cổ vũ nhiệt liệt.

Đương nhiên, những khán giả ``bất đắc dĩ'' - tức là các cô gái \hll\ và cả Kuon cùng
A-chan - cũng trầm trồ ngạc nhiên với kỹ thuật múa lân điệu nghệ của nàng Oni này.

``Không ngờ cô ấy lại làm tốt như vậy\dots'' Aki trầm trồ, ánh mắt không rời khỏi Ayame.
``Chắc tại nhỏ người nên dễ di chuyển hơn,'' Matsuri vừa nói vừa cười tươi.

``Quả không ngoa khi nói rằng em ấy có kỹ thuật tốt thật,'' Kuon khá bất ngờ với màn
trình diễn của cô ấy, dù cho tỏ vẻ khá cảnh giác: ``Cơ mà anh nghĩ em ấy sẽ lại làm gì dở
hơi thôi-''

Vừa dứt lời, màn biểu diễn kết thúc trong tiếng cười bùng nổ khi Ayame, trong lúc cố
gắng thực hiện cú nhào lộn, đã khiến cả con lân ``ngã chổng vó'' xuống đất.

Dù vậy, nàng Quỷ sừng vẫn cười tươi rói, không hề có chút ngại ngùng. Cô thốt lên, hùng
hồn tuyên bố: ``Cái này\dots\ gọi là nghệ thuật gây bất ngờ!'', nhận về tràng pháo tay vang
dội từ khán giả.

Sau khi hoàn thành màn trình diễn đầu tiên đầy bất ngờ, nàng Quỷ sừng bỗng lóe lên một
ý tưởng táo bạo khác. Nhìn thấy một đội múa lân khác đang chuẩn bị, cô không kiềm chế
được sự phấn khích.

Cô chạy đến gần đội đó, nở nụ cười rạng rỡ: ``{\color{red} Cho tôi và vài người bạn của tôi thử được
  không?}''

Một thành viên trong đội múa hơi ngạc nhiên, nhưng rồi hỏi với vẻ nghi ngờ: ``{\color{red} Vậy đồng
  đội của cô đâu?}'' Không để mất thời gian, Ayame nhanh chóng kéo tay Subaru, Mio, và
Fubuki đứng vào hàng, thốt lên: ``{\color{red} Có liền!}'' khiến cho anh chàng đội múa đổ mồ hôi hột.

Bản thân nàng Vịt Subaru cũng tỏ vẻ bất ngờ và có đôi chút không đồng tình với nàng
Quỷ sừng nhiệt tình kia: ``Này! Tôi có đồng ý tham gia chuyện này đâu!'' Cô vùng vàng,
cố gắng thoát khỏi bàn tay nhỏ nhưng chắc của Ayame.

``Tôi còn chẳng biết nó là gì mà\dots'' Mio lo ngại, lên tiếng phản đối, đôi mắt đầy lo lắng.

Ngược lại, Fubuki lại bật cười, trông có vẻ háo hức hơn hẳn: ``Cứ thử đi xem nào, biết
đâu lại vui thì sao!''

Một người trong nhóm múa đứng gần đó nhìn cảnh tượng kỳ lạ trước mắt, vẻ mặt không
giấu được sự ái ngại: ``{\color{red} Trông hai bạn tóc ngắn với bạn tai chó hoang mang như vậy\dots\ Liệu họ có làm nên cơm cháo gì không đây\dots}''

Người đồng đội bên cạnh chỉ nhún vai: ``{\color{red} Thôi kệ, cứ để họ thử xem sao.}''

Người đó chỉ biết thở dài, rồi miễn cưỡng trỏ tay vào bộ đồ lân đã chuẩn bị sẵn: ``{\color{red} Mọi
  người cứ việc.}'' Ayame lập tức chạy đến, ánh mắt sáng lên đầy quyết tâm.

Dường như với bốn người trong bộ tứ FAMS\footnote{Về tên gọi của nhóm này, xem lại trong {\flaregothic ÉPISODE 21}, quyển số Hai.} là chưa đủ, cô nhanh chóng kéo tay Miko,
người đang đứng phía sau đám đông, vào cuộc. ``Miko-chan, em nhờ chị làm chỉ huy nha,
dazo!''

Nàng Vu nữ phản ứng lại với yêu cầu bất ngờ của cô nàng, giọng lúng túng rõ rệt: ``Ơ,
nhưng\dots\ chị có biết gì về mấy thứ này đâu ne!?''

``Chị cứ làm thử đi!'' Ayame đáp chắc nịch, giọng nói như thể khẳng định rằng đối phương
không được phép từ chối. Hết cách, Miko chỉ biết thở dài cam chịu, rồi đi theo cả nhóm
đến gần bộ đồ lân.

Với năm người lính mới, họ bắt đầu chia nhau vai trò trong con lân: Subaru và Mio làm
chân sau, Fubuki và Ayame phụ trách phần đầu, còn Miko thì đứng ngoài với vai trò chỉ
huy, đảm nhận nhiệm vụ phối hợp mọi người.

\ct

Nhóm FAMS đứng ở vị trí trung tâm của đám đông, sẵn sàng cho màn trình diễn độc đáo.
Nàng Vu nữ Miko-chi, giờ đã thay bộ trang phục vu nữ năm mới\footnote{Xem hình d, trang \pageref{settwo}.}, đứng phía trước với
cây gậy trừ tà trong tay. Cô vung nhẹ cây gậy, tạo dáng như thể mình là nhạc trưởng của
một dàn nhạc thính phòng.

Ở phía bên cạnh, Korone và Okayu, bị Ayame lôi kéo làm ``tay trống'' bất đắc dĩ, đang
chuẩn bị với vẻ mặt vừa buồn cười vừa khó hiểu.

Nàng Khuyển lẩm bẩm: ``Cậu ấy bảo mình đánh trống như chó sủa\dots\ Nhưng mà trống thì
làm sao sủa được chứ!?'' Nàng Miêu cười nhẹ, vỗ nhẹ vào trống: ``Ít ra chúng ta không
phải chui vào bộ đồ lân đó, nhỉ\~{}?''

``Còn tùy xem việc này có kéo dài đến bao lâu nữa\dots'' cô bạn thở dài, nhưng rồi cũng cố
nở nụ cười. ``Thôi thì, vì mọi người, mình sẽ cố hết sức!''

Nàng Oni quay đầu lại, vẫy tay động viên hai người: ``Cố lên nhé, hai đồng chí trống chiến!
Tôi đặt kỳ vọng vào mọi người đấy!''

Cặp Chó-Mèo nhìn nhau rồi bật cười, lắc đầu: ``Đúng là không thể thắng được sự nhiệt
tình của cô ấy thật á.''

Trong khi đó, các thành viên còn lại đứng trong đám đông khán giả, quan sát màn chuẩn
bị. Kuon nheo mắt nhìn nhóm FAMS và thì thầm: ``\textit{Sao mình thấy có gì đó sai sai nhỉ\dots}''
nhưng rồi cũng tặc lưỡi cho qua, đồng thời cũng hiếu kỳ: ``Không biết lần này họ sẽ có
bài gì cho chúng ta xem đây\dots''

Miko hít một hơi sâu, sau đó vung gậy lên và ra chỉ thị bắt đầu. ``Lên đi nào, các em!''

Cặp Chó Mèo liền gõ nhịp đầu tiên, từng hồi trống vang lên đầy mạnh mẽ và đều đặn.
Con lân của nhóm FAMS bắt đầu chuyển động, từng phần được phối hợp nhịp nhàng.

Phần đầu lân do Ayame điều khiển lắc qua lắc lại một cách đầy năng lượng. ``Không biết
lắc lắc như thế này có đúng không nhỉ\dots'' cô nàng vừa tự hỏi, vừa cố gắng làm theo những
gì cô nhớ từ màn trình diễn trước đó.

Phần thân lân ở phía sau, do Mio và Subaru đảm nhận, cũng bắt đầu chuyển động. Nàng
Sói đen nghiêng người nhẹ, thử nghiệm: ``\textit{Mình nghĩ chắc là nên lắc như thế này\dots}'' theo
sau là nàng Vịt đang nhìn vai, rồi cũng bắt chước: ``\textit{Chắc là làm thử như Ayame-chan xem
  sao.}''

Fubuki, đảm nhận phần đuôi, không chịu thua kém, vẫy mạnh bộ lông lân trong tay, hú
hét: ``\textbf{Woo-hoo! Trông mình ngầu không này!}''

Mio và Subaru đồng thanh quát lên, làm nàng Cáo trắng giật mình: ``\textbf{Trật tự đi giùm chúng tôi cái!!}''

Dưới sự dẫn dắt của Miko, nhóm bắt đầu di chuyển vòng quanh khán giả. Tiếng trống
dồn dập hơn khi con lân đi vào nhịp điệu.

Miko hất nhẹ cây gậy trừ tà lên trời. Thấy thế, nhóm FAMS đồng loạt nhảy lên theo. Gậy
vung thành hình cung, họ lập tức di chuyển giống như một đoàn rồng rắn lên mây.

Khán giả vỗ tay reo hò, tiếng cổ vũ càng làm không khí trở nên sôi động. Màn trình diễn
vẫn chưa kết thúc, nhưng những nụ cười rạng rỡ và sự hào hứng của khán giả đã cho thấy
họ thực sự bị cuốn hút vào màn múa lân ``phiên bản \hll\ sáng tạo'' này.

Bầu không khí lễ hội càng thêm rộn ràng khi nhóm bộ tứ ``giọng trầm'' tiếp tục phần trình
diễn độc đáo của mình bằng những động tác bán chuyên nghiệp. Đám đông khán giả đứng
vây quanh, vừa tò mò vừa thích thú trước ``con lân'' lạ lẫm và tiếng trống vang dội.

Một người trong đám đông chậc lưỡi, chỉ tay: ``{\color{red} Con lân này trông kỳ lạ ghê!}''

Một người khác bật cười: ``{\color{red} Trông nó múa buồn cười quá!}''

Còn một người đứng gần cặp Chó Mèo lại thì thầm: ``{\color{red} Không biết cặp đôi đánh trống kia
  có phải là dân cosplay không nhỉ\dots}''

Tiếng trống càng lúc càng dồn dập, tạo nên nhịp điệu sôi động. Trong khi đó, Miko, đứng
ở vị trí chỉ huy phía trước, hăng say vung gậy trừ tà. Nhưng rõ ràng, sự tự tin của cô nàng
không đi kèm với kỹ thuật thuần thục. Cô nàng lẩm bẩm, tay vẫn cầm cây gậy không
ngừng chỉ tứ tung: ``\textit{Như thế này\dots\ đúng không nhỉ, ne?}''

Từ phía khán giả, Kuon nhếch miệng cười, nói nhỏ với A-chan: ``{Em ấy trông có vẻ làm
  tốt đấy\dots}''

Nàng Đeo kính khẽ gật đầu: ``Công nhận. Nhưng mà\dots\ trông hơi lạ.''

Không những vậy, trong lúc Miko hăng hái chỉ huy, cô còn vô tình tạo ra một số tư thế
đặc biệt. Những động tác này chẳng những thu hút ánh mắt ngạc nhiên của khán giả mà
còn khiến các đội múa lân khác phải há hốc mồm kinh ngạc.

Dường như không để ý tới đội múa của mình, Miko xoay gậy vòng quanh, tạo nên một
vòng cung lớn, rồi hét lên: ``\textbf{Tiến lên! Mọi người nhảy đi nào!}''

Như thể bị lây năng lượng của nàng Vu nữ, Ayame hô lớn: ``\textbf{A-lê-hấp!}'' Con lân đột nhiên
nhảy vọt lên theo hiệu lệnh của Miko, làm cả đám đông phải bật cười, đồng thời cổ vũ
một cách nhiệt liệt.

Subaru ở phía sau, không giấu được vẻ ngỡ ngàng trước hành động bất ngờ của cô bạn,
hét toáng lên: ``\textbf{Này- Bà làm cái gì vậy!?}''

Tiếng trống lại vang lên dồn dập hơn, như tiếp thêm lửa cho sự náo nhiệt của màn trình
diễn. Nhưng những động tác bất ngờ của nhóm FAMS khiến cả khán giả và những nhóm
đang trình diễn khác đều không thể rời mắt.

Những tiếng reo hò của khán giả càng lúc càng lớn, như tiếp thêm sự hăng hái cho nhóm
FAMS. Cả bốn người họ phối hợp một cách lạ kỳ, thực hiện màn xoay vòng trên không
trung hết lần này đến lần khác, dù không ai thực sự biết làm thế nào mà mọi thứ vẫn\dots\ không bị sụp đổ.

Miko vẫn đang vung vẩy cây gậy một cách tứ tung, đến mức chính cô nàng còn không
hiểu mình đang làm gì: ``À-rế, ne? Mọi người đang làm cái gì thế này?''

Dù thế, cô vẫn tiếp tục chỉ huy đội múa lân của mình, như thể luôn tự nhủ rằng màn trình
diễn phải tiếp tục dù nó đang xảy ra một cách loạn xạ và lung tung đến cỡ nào.

Bên trong bộ đồ lân, Mio bắt đầu không thể theo kịp được nàng Oni đang hăng say kia.
Cô hét lên, giọng đầy hoảng loạn: ``\textbf{Chúng ta có phải là diễn viên xiếc đâu! Miko-chan,
  chị đang làm cái hú gì vậy!?}''

Trái ngược với sự sợ hãi của nàng Sói, Fubuki lại cười phá lên, tỏ vẻ thích thú: ``\textbf{Vui quá
  xá là vui! Yahoo! Lên nào!}''

Sau đó, cả nhóm tiếp tục nhảy và xoay thêm vài vòng trên không trung. Nàng Quỷ sừng
phấn khích hô vang: ``\textbf{Weee!}'' theo sau đó là nàng Cáo trắng vẫn đang cười lớn, như thể
họ đang đi trên tàu lượn siêu tốc.

Mio và Subaru thì dường như bắt đầu hoang mang vì cặp đôi nhiệt tình này. Nàng Sói
đen trong khi cố gắng giữ thăng bằng, run rẩy nói: ``C-Chóng mặt quá\dots''

Đỉnh điểm của màn trình diễn điên rồ là khi họ quyết định xoay toàn thân con lân trên
không trung, khiến Subaru tái mặt, hét toáng lên: ``\textbf{M-Miko-chan, dừng lại đi!}''

Nhưng nàng Quỷ sừng, vẫn giữ tinh thần hăng hái không gì cản nổi, hét lớn: ``\textbf{Nào ta cùng bay lên nào, hỡi chị em FAMS ơi!}''

``Á, đừng-'' Mio định kêu lên, nhưng rồi bị tiếng hô hào của Ayame chặn họng. Nàng Sói
đen, cùng với Fubuki và Subaru, cùng nhau nhảy lên trên, tạo thành một ``sóng lân''. Hàng
tiếng người cổ vũ gầm rú lên, khiến cho bầu không khí xung quanh họ càng náo nhiệt hơn
bao giờ hết.

Ở phía sau, Korone và Okayu dường như đã bị cuốn theo không khí náo nhiệt. Tiếng
trống mà họ đánh vang lên càng lúc càng dồn dập, thúc giục khán giả xung quanh đồng
thanh hô lớn: ``\textbf{Lên! Lên! Lên!}''

Miko, ở vị trí chỉ huy, lắc lắc gậy trừ tà như thể đang dẫn dắt cả một buổi biểu diễn chuyên
nghiệp. Nhưng trên khuôn mặt cô hiện rõ vẻ vừa hoang mang vì cô vẫn không biết cô
đang làm trò gì, lại vừa tỏ vẻ quyết tâm khi buổi biểu diễn của nhóm trong con mắt của
khán giả đang rất tích cực.

``\textit{Không còn cách nào khác rồi\dots\ Đã đâm lao thì phải theo lao rồi, ne!} \textbf{Lên nào, mấy đứa!}''

Tiếng hò reo càng lúc càng mãnh liệt, bầu không khí náo nhiệt bao trùm toàn bộ khu vực.
Màn trình diễn đã vượt ra khỏi tầm kiểm soát ban đầu, nhưng sự phấn khích của cả nhóm
và khán giả đã biến tất cả trở thành một cảnh tượng đáng nhớ.

Màn trình diễn tưởng chừng đã lên đến đỉnh cao, nhưng không ngờ, bất ngờ khác lại ập
đến. Từ dưới đất, hàng loạt cột tre có gắn bệ đứng bất ngờ mọc lên như nấm, buộc cả
nhóm FAMS phải điều khiển con lân đáp xuống.

Mio hét lớn, giọng đầy tuyệt vọng: ``\textbf{Cái gì nữa đây hả trời!!}''

Ayame hoàn toàn phớt lờ sự lo lắng của đồng đội, hứng khởi kêu lên: ``Mọi người cùng
biểu diễn nào!''

Họ không còn lựa chọn nào khác ngoài việc tiếp tục màn trình diễn, lần này là trên những
cây tre chênh vênh. Con lân do bốn người điều khiển bắt đầu thực hiện những động tác
vô cùng điêu luyện, từ nhảy cóc, lộn nhào, cho đến đu dây.

Điều đáng kinh ngạc là hình dáng của con lân vẫn được giữ nguyên vẹn, như thể họ là
những chuyên gia thực thụ.

Fubuki hét vang, giọng đầy phấn khích: ``\textbf{Woo-hoo! Đúng là đỉnh của chóp luôn!}'' Nàng
Quỷ sừng cười rạng rỡ, tiếp thêm năng lượng cho cả nhóm: ``\textbf{Đúng rồi! Cứ thế mà làm!}''

Trái ngược hoàn toàn, Subaru bấu chặt lấy phần khung bên trong con lân, hét lên trong
cơn hoảng loạn: ``\textbf{Tôi không muốn làm diễn viên xiếc đâu!}''

Mio cũng run rẩy không kém. Mồ hôi cô chảy ròng, giọng lạc hẳn đi: ``Làm ơn! Ai đó
hãy dừng cái này lại đi!''

Ở phía ngoài, Kuon và A-chan đứng quan sát với ánh mắt pha lẫn giữa kinh ngạc và thích
thú. Cậu Tổng vỗ tay bôm bốp, dường như không để ý tới hai tâm trạng trái ngược bên trong con lân đó: ``Quả đúng là diễn viên xiếc thật.''

Nàng Đeo kính gật gù, mắt vẫn dán chặt vào màn trình diễn: ``Tôi cũng thấy ngạc nhiên
với trình độ xiếc của họ đấy.''
Bên trong con lân, Ayame và Fubuki vẫn vui vẻ hét lớn, hoàn toàn phớt lờ sự tái mặt của
hai người còn lại.

\begin{dialogue}
  \speak{Ayame} \textbf{Đu dây nào! Bay lên trời cao!}
  \speak{Fubuki} \textbf{Thêm một vòng nữa! Tuyệt quá!}
  \speak{Ayame} \textbf{Thôi nào, Subaru, Mio-chan, hai người không cảm thấy hào hứng sao!?}
\end{dialogue}

Subaru tiếp tục la lên, giọng cô khản đặc dần vì la hét: ``\textbf{Làm ơn cho tôi xuống\dots\ Tôi
  không chịu nổi nữa rồi! M-Miko-chan\dots\ dừng lại!}''

Mio dù đã cố gắng giữ bình tĩnh, cuối cùng cũng hét lên: ``\textbf{Ayame! Dừng lại đi!
  Tôi sắp rớt mất rồi!}''

Nhưng bầu không khí náo nhiệt và tiếng reo hò cổ vũ của khán giả dường như đã cuốn tất
cả vào một màn trình diễn mà không ai ngờ tới. Những động tác điêu luyện trên những
cây tre mỏng manh càng làm khán giả phấn khích hơn bao giờ hết.

Khi những động tác trên cây tre đang tiến dần đến những điệu cuối cùng, Miko bất chợt
nhảy cẫng lên, giơ cao cây gậy trừ tà như một nhạc trưởng đầy khí thế. ``\textbf{Đây là đòn cuối
  cùng đây, ne!}''

Mọi người dừng lại, tập trung theo dõi chỉ dẫn của nàng Vu nữ. Nhưng bất ngờ thay,
từ trong con lân, nàng Quỷ sừng bất ngờ lấy ra một chai tương ớt đỏ rực, khiến Mio và
Subaru đồng loạt hét lên: ``\textbf{B-Bà kiếm cái này ở đâu ra vậy!?}''

Ayame mỉm cười đầy tự hào: ``Từ khi nãy á!'' Subaru trố mắt, cảm giác bất an ngày càng
lớn: ``N-Này, đừng nói với tôi là bà định-''

Chưa kịp dứt lời, Ayame cầm chai tương ớt tu một hơi sạch bóng. Hai cô nàng xanh mặt
đồng loạt thốt lên, tiếng hét vang vọng cả khu vực: ``\textbf{Đừng mà!!}''

Ngay sau đó, nàng Oni nghiêng đầu ra phía trước, phun ra một ngọn lửa sáng rực, tạo
thành một vòng cung lửa tuyệt đẹp giữa không trung. Khán giả bên dưới bùng nổ tiếng
reo hò cổ vũ.

Kuon đứng nhìn, khóe miệng khẽ nhếch lên, bình thản nhận xét: ``Giờ mới đúng chuẩn là
`múa lân phiên bản \textit{hologra}' thật này.''

A-chan chỉ biết thở dài, nhưng vẫn không giấu nổi nụ cười.

Nhóm những cô gái còn lại có những phản ứng rất riêng, nhưng nhìn chung là họ cảm
thấy ngưỡng mộ với màn trình diễn tuyệt hảo này.

Choco, mắt sáng rực như thấy thứ gì đó rất hữu ích: ``Kỹ thuật này\dots\ làm cho con lân này
biến thành con rồng rồi á!''

Aki lắc đầu, nhưng không giấu được sự ấn tượng: ``Đúng là một cách biểu diễn nghệ thuật
độc đáo\dots\ Hình như mình cũng muốn thử một lần.''

Pekora thì ôm bụng cười không ngớt, chỉ tay về phía con lân đang phun lửa kia: ``Ha \uparrow ha \downarrow ha \uparrow ha \downarrow Không ngờ Ayame-senpai lại có tài lẻ kiểu này luôn á, peko!''

Mel che miệng cười khúc khích: ``Nóng bỏng quá\dots\ Nhưng mà có vẻ hơi nguy hiểm thì
phải?''

Khán giả vẫn tiếp tục vỗ tay rầm rộ, trong khi Ayame đứng giữa, thỏa mãn vẫy tay chào
đám đông, như thể ngọn lửa vừa rồi là ``chiêu kết thúc tuyệt đối'' của cả màn trình diễn.

Khi tiếng trống của Korone và Okayu dần kết thúc, cả hai đồng loạt gật đầu với nhau,
đếm nhịp một cách đầy hào hứng.

\begin{dialogue}
  \speak{Korone} Se-no (Hai, ba)!
  \speak{Okayu} Đánh nào~
\end{dialogue}

Cả hai cùng lúc giáng dùi vào chiếc cồng chiêng lớn được dựng sẵn phía sau.

\textbf{*ĐOÀNG!*}

Một âm thanh trầm vang rền khắp lễ hội, khiến cả đám đông theo phản xạ lùi lại một bước,
nhưng ngay lập tức họ vỡ òa trong những tràng pháo tay rôm rả.

Những người xung quanh vẫn không ngừng hò reo cổ vũ, ngay cả khi màn trình diễn đã
kết thúc. Tiếng cười và tiếng nói pha lẫn vào nhau tạo thành không khí đầy náo nhiệt.

Tất cả các thành viên \hll\ còn lại, đứng xem từ xa, vẫn chưa nói được lời nào, hoàn
toàn bị ấn tượng bởi màn trình diễn.

Rushia tròn xoe mắt: ``Tuyệt thật đó-nanodesu\dots'' Sora dịu dàng, nhưng giọng vẫn không
giấu nổi sự ngạc nhiên: ``Đây là lần đầu tiên chị được xem màn múa lân tuyệt vời như
vậy\dots''

Roboco gật gù, giọng điệu pha chút ngập ngừng: ``Chị không nghĩ chị có thể làm được
như họ đâu\dots''

Aki mỉm cười, đôi mắt lấp lánh, như thể cô nàng vừa xem một màn trình diễn bùng nổ
nhất trong cuộc đời mình: ``Đúng là đã mắt thật!''

Pekora, như thường lệ, khoanh tay với nụ cười đầy thách thức: ``Tôi có thể làm được tốt
hơn đó, peko.''

Shion lắc đầu, phồng má phụng phịu: ``Đáng ra người điều khiển phải là em chứ!''

Noel, nhìn chằm chằm vào con lân với ánh mắt mơ mộng, buột miệng: ``Trông giống như
hiệp sĩ đi đánh rồng vậy!''

Matsuri, không giấu nổi vẻ phấn khích, nhảy cẫng lên: ``Nhìn đúng là vui quá trời vui!''

Bản thân Kuon và A-chan cũng đều mãn nguyện với màn trình diễn tuyệt đỉnh của FAMSGAMERS.
Họ nhìn nhau, không nói một lời nào, rồi hòa cùng đám đông trao cho họ một tràng vỗ tay giòn giã.

Ở phía xa, Ayame, Fubuki, Mio, Subaru, Miko, Korone, và Okayu cùng cúi chào đám
đông một cách đồng bộ, khép lại màn trình diễn đầy màu sắc và không kém phần điên rồ.
Tiếng hò reo, tiếng vỗ tay, và tiếng hô vang vẫn tiếp tục cho đến khi bảy người rời khỏi
sân khấu.

\ct

Sau màn biểu diễn rực lửa, nhóm múa lân tản ra nghỉ ngơi. Ayame, khuôn mặt đỏ ửng vì
sức nóng từ tương ớt, phải ngồi một góc, tay ôm túi đá lạnh lên má mà vẫn không ngừng
kêu đau.

``Hức\dots\ đau quá\dots\ không ngờ nó lại cay thế này\dots'' nàng Quỷ sừng ôm túi đá lạnh chườm,
khẽ run rẩy tới mức rơi nước mắt.

Một thành viên của đội múa lân khác đi ngang qua, tò mò hỏi thăm: ``{\color{red} Cô gái, sao lại đến nỗi thế này?}''

Cô đáp, mặt méo xệch vì đau: ``{\color{red} Tôi chỉ\dots\ thử chút tương ớt thôi mà\dots\ ai ngờ\dots\ Ái cha!}''

Nghe cô nàng đang mếu máo kia, người đó bật cười, đồng thời vỗ vai cô nàng động viên:
``{\color{red} Buổi trình diễn của mấy cô tốt thật đấy. Nhưng mà lần sau, biết ớt cay như thế thì chừa
  nó ra nhé.}''

Tất cả những người xung quanh nghe vậy cũng đều bật cười với hình dáng tội nghiệp của
cô nàng Oni hăng máu ấy.

Mio và Subaru đứng gần đó, khoanh tay thở dài đầy chua chát:

\begin{dialogue}
  \speak{Mio} Bọn tôi đã bảo đây là ý tồi rồi mà\dots
  \speak{Subaru} Thế mà vẫn không nghe\dots
\end{dialogue}

Choco nhẹ nhàng bước tới, cười duyên dáng: ``Thôi nào, mọi người. Có màn trình diễn
nào mà không cần chút mạo hiểm đâu. Ayame-sama đã làm rất tốt mà.''

Nàng Quỷ sừng vội nắm tay nàng Y tá, mắt long lanh: ``Choco-sensei\dots\ Cô là thiên thần
của em đó!!''

Nàng Vịt ngao ngán đáp, cảm nhận thấy ánh sáng của thiên sứ đang tỏa lên quanh người
cô: ``Cái này nghe quen quen đấy\dots'' Nàng Sói đen thì ngán ngẩm không nói lấy một câu
nào.

Ở một góc khác, bốn người vừa hoàn thành phần biểu diễn cùng nhau trò chuyện vui vẻ.
Nàng Vu nữ vỗ tay bôm bốp với cặp đôi Chó-Mèo và Cáo trắng: ``Mấy đứa đánh trống
dữ quá ne~''

Okayu nhún vai, điểm tĩnh đáp: ``Bọn em có chơi Taiko (no) Tatsujin (太鼓の達人) trước đó rồi mà, Miko-chan\~{}''

Korone ôm lấy vai cô bạn, hớn hở thêm vào: ``Còn em thì được đánh chung với Okayu\~{} Thật là vui quá đi!''

``Cũng nhờ công Miko-chan đã điều khiển cả buổi biểu diễn này mừ\~{}'' nàng Miêu nói,
khen lấy nàng Vu nữ, khiến cho đối phương đỏ mặt ngượng ngùng.

``Thật sự thì chị cũng không biết chị vừa làm cái gì nữa ne\dots\ Nhưng mà, cảm ơn em.''

Fubuki bật cười, ánh mắt lấp lánh sự cảm kích: ``Chị cũng không nghĩ là mấy đứa đỉnh
đến thế đấy! Cảm ơn nha!''

Từ phía xa, A-chan và Kuon quan sát nhóm bảy người đang đùa giỡn, trò chuyện rôm rả
với nhau.

Nàng Đeo kính khoanh tay, khẽ mỉm cười: ``Nhìn mấy đứa kia trông cứ như một gia đình
nhỏ ấy nhỉ.''

Cậu Tổng gật đầu, đôi mắt không rời khỏi khung cảnh náo nhiệt trước mặt: ``Chị nói cũng
phải. Nhưng mà gia đình này\dots\ có lẽ là hơi ồn ào đấy.''

``Ồn ào thật. Nhưng cũng đáng yêu mà, đúng không?''

``Dạ\dots\ cũng phải ạ. Dù gì hôm nay cũng là ngày cuối cùng họ sẽ đón Tết ở đây rồi, nên
chắc chúng ta cũng phải vui chơi một chút nhỉ?''

``Đồng ý.''

Tiếng cười đùa của nhóm bảy người tiếp tục vang vọng trong không khí mát mẻ của buổi
tối lễ hội, khép lại một ngày không thể nào quên.

Buổi tối ngày Tết cuối cùng, các thành viên của \hll\ cũng tham gia trải nghiệm các
hoạt động đầu xuân năm mới ngay tại đât tổ của người Etsunan.

Chẳng hạn như\dots

\il

Matsuri đứng giữa một gian lều nhỏ bên cạnh những ông đồ thực thụ, xung quanh là các
bộ dụng cụ viết thư pháp bày biện gọn gàng. Mặc chiếc áo dài màu vàng tươi và đội khăn
đóng, cô nàng trông không khác gì một ông đồ ``phiên bản thần tượng'' cả.

Tuy nhiên, khuôn mặt cô lại phảng phất sự bối rối, khi đã ngồi suốt khoảng nửa tiếng,
nhưng chưa ai ghé vào quầy của cô cả. ``{\color{red} Ơ\dots\ sao chẳng có ai lại xin chữ mình nhỉ? Kỳ lạ
  quá\dots}''

Bên cạnh quầy của cô, một ông đồ già, vừa cầm bút lông vừa nhìn nàng Cổ động, lắc đầu
cười nhẹ: ``{\color{red} Cô gái trẻ, nghề này không phải là trò đùa đâu. Chữ nghĩa phải có tâm và tầm,
  chứ không phải muốn cho là cho.}''

Matsuri chớp chớp mắt, cố giữ nụ cười: ``{\color{red} Ơ, cháu tưởng việc xin chữ giống như ký tặng
  ấy ạ. Là thần tượng \hll\ như cháu mà không có người nào nhận ra sao ạ?}''

Ông đồ già bật cười, đáp lại bằng giọng điềm đạm: ``{\color{red} \hll\ là cái gì nhỉ? Ta chưa nghe
  đến bao giờ.}''

Nghe vậy, Matsuri tròn mắt, biểu cảm như không tin vào tai mình. Cô quay lại nhìn Kuon,
người đang đứng gần đó cùng A-chan, để tìm sự xác nhận: ``Kuon-senpai, ông ấy nói thế
là sao!? \hll\ mà cũng không biết ư!?''

Cậu Tổng nhún vai, cố nén cười: ``Matsuri-chan\dots\ ở đây không phải Nhật Bản đâu, đây
là Etsunan mà. Có phải ai cũng biết công ty chúng ta đâu chứ.''

A-chan cầm một tờ giấy đỏ khổ lớn, quay sang cười khúc khích: ``Nhưng mà nếu Natsuirosan
muốn, sao không thử viết vài chữ xem? Biết đâu có người thấy hứng thú thì sao?''

Nghe lời động viên từ cô, nàng Lễ hội lấy lại tinh thần. Cô xắn tay áo dài, nắm chặt cây
bút lông trong tay: ``Được rồi, để xem tài năng của Matsuri-sama này sẽ khiến họ trầm
trồ thế nào!''

Nói là làm, cô bắt tay viết những chữ lên tấm giấy đỏ một cách bạo dạn, hừng hực khí thế.

``Mình không nhớ là thầy đồ sẽ viết chữ đầy mạnh bạo như thế\dots'' Kuon nhìn cô nàng viết
xoèn xoẹt chữ thầy đồ như thể cô nàng đang viết tặng chữ ký.

\begin{center}
  \uline{Vài phút sau.}
\end{center}
Một số người qua đường nhìn các chữ mà cô nàng treo trên quầy, có người thì tỏ vẻ tò
mò, có người thì ngơ ngác không hiểu, có người thì lại thấy hứng thú.

``{\color{red} Chữ này\dots\ là chữ gì thế nhỉ? Trông lạ quá\dots}''

``{\color{red} Hán Nôm không ra Hán Nôm, Quốc Ngữ không ra Quốc Ngữ nữa\dots\ Đây có phải kiểu
  thư pháp hiện đại không?}''

``{\color{red} Tôi tưởng idol thì chỉ hát với nhảy thôi, không ngờ lại biết cả viết thư pháp\dots}''

Matsuri nghe loáng thoáng những lời bàn tán, nụ cười trên môi dần méo mó. Cô quay lại
nhìn Kuon, mắt như muốn khóc: ``K-Kuon-senpai\dots\ cứu em với! Họ có vẻ không hiểu
nghệ thuật của em kìa!''

Cậu Tổng nghe cô nàng than vãn như thế, vừa cười vừa xua tay: ``Thôi thì, ít nhất thì em
cứ coi như đây là buổi thử nghiệm đi. Biết đâu lần sau sẽ tốt hơn.''

A-chan đặt một tay lên vai của cô nàng đang mếu máo kia, không nhịn được cười: ``Dù
sao thì cô cũng đã thử sức hết mình rồi mà. Đó cũng là một cách để lan tỏa văn hóa
\hll\ mà, đúng không?''

Nghe những lời động viên từ hai người, nàng Lễ hội lưỡng lự một lúc lâu, rồi mỉm cười
gượng gạo: ``Hai người nói cũng phải\dots\ Chắc lần sau em sẽ phải học hỏi kỹ lưỡng hơn
nhỉ?''

\ct

Ở một diễn biến khác, Choco bước chậm rãi qua các gian hàng, chiếc áo dài màu đỏ thêu
hoa mai làm tôn lên vẻ đẹp thanh lịch của cô. Ánh mắt của cô nàng long lanh khi nhìn
thấy những nét bút thư pháp uyển chuyển trên những tấm giấy đỏ. Cô quyết định tiến đến
một ông đồ đang mải mê vẽ từng đường nét trên tờ giấy trước mặt.

``{\color{red} Cháu chào ông ạ!}'' cô nàng cúi chào ông đồ đầy lễ phép. ``{\color{red} Cháu muốn xin một chữ để
  mang về làm kỷ niệm được không ạ?}''

Trước mặt cô, một cụ già với mái tóc bạc phơ và đeo cặp kính lão. Cụ ngẩng đầu lên mỉm
cười: ``{\color{red} Chào cô gái trẻ. Muốn xin chữ gì nào? Chữ `Phúc', `Lộc', hay `Thọ'?}''

Nàng Y tá nghiêng đầu, đôi mắt mơ màng: ``{\color{red} Chữ nào mang ý nghĩa tình yêu hay sự ngọt
  ngào ạ. Vì cháu là cô giáo dạy môn\dots\ yêu thương mà.}''

Kuon nhìn cô nàng nháy mắt và nói với giọng khiêu gợi, cậu hắng giọng một tiếng:
``Choco-sensei, đây không phải chỗ để nói linh tinh đâu.''

Ông đồ bật cười, đôi tay thoăn thoắt viết một chữ lên giấy đỏ. Sau vài nét bút dứt khoát,
ông đặt bút xuống, đưa tấm giấy cho Choco: ``{\color{red} Vậy ông sẽ cho cháu chữ `Tâm' nhé. Nó
  mang ý nghĩa là trái tim, lòng thành, và tình yêu thương. Rất hợp với cô gái trẻ như cô.}''

Cô nàng vui vẻ nhận lấy tấm giấy đỏ mà ông đồ vừa viết. Nhưng ngay khi nhìn vào chữ
viết, cô khựng lại. Khuôn mặt đang tươi rói bỗng đờ ra, rồi cô lẩm bẩm: ``{\color{red} Ơ\dots\ Chữ này
  là chữ gì vậy nhỉ? Sao cháu không đọc được?}''

Ông đồ bật cười, nhẹ nhàng giải thích: ``{\color{red} Đây là chữ Hán Nôm, cô gái. Chữ này thể hiện
  tinh thần truyền thống và nét đẹp văn hóa của người xưa.}''

Nàng Y tá nghe ông đồ giải thích lại càng bối rối hơn. Cô xoay xoay tấm giấy, nhìn nét
chữ vừa lạ vừa quen đó, như cố gắng giải mã từng nét chữ: ``{\color{red} Nhưng mà\dots\ cháu chỉ biết
  đọc chữ Kanji thôi, chứ chữ Hán Nôm thì cháu chịu\dots}''

Nghe thấy vậy, Kuon và A-chan, đang đứng gần đó, không nhịn được cười. Cậu Tổng
tiến lại gần, cầm lấy tấm giấy: ``Sensei à, đó là chữ `tâm' đó. Đúng như ông ấy nói, nó có
nghĩa là trái tim, nhưng cũng ám chỉ lòng nhân hậu và sự chân thành nữa. Em không cần
phải đọc được, chỉ cần cảm nhận ý nghĩa là được rồi.''

Choco chớp mắt, khuôn mặt dần giãn ra khi được cậu ``giải ngố.'' Cô mỉm cười dịu dàng:
``Ra vậy\dots\ Đúng là văn hóa truyền thống sâu sắc thật. {\color{red} Cảm ơn ông đã giải thích cho
    cháu!}''

Ông đồ cười hài lòng: ``{\color{red} Không có gì. Văn hóa là để chia sẻ, cô gái ạ.}''

Nàng Đeo kính cười nhẹ: ``Yuzuki-san, có khi cô nên giữ tấm giấy này làm kỷ niệm. Ai
biết đâu sau này chị sẽ nhớ được một chút về chữ Hán Nôm thì sao.''

Nàng Y tá gật đầu, áp tấm giấy vào ngực, ánh mắt đầy trân trọng: ``Cảm ơn chị ạ. Chắc
có lẽ mỗi lần em nhìn vào nó, em sẽ nhớ tới chuyến đi này.''

\ct

Tại góc nhỏ phía trước sân đền, Miko đã dựng hẳn một quầy rút quẻ với bảng hiệu viết
tay ghi dòng chữ bằng tiếng Etsunan: ``\uline{Rút quẻ lấy may! Đảm bảo chính xác 100\%, nếu
  không thì\dots\ cố gắng lần sau ne\~{}!}'' với chữ ``ne'' được viết là ``にぇ'' như thể đây là câu
khẩu hiệu của mình.

Cô nàng vẫn mặc bộ đồ vu nữ đặc trưng, tay cầm gậy trừ tà vung vẩy như đang chỉ huy
đội lân như hồi nãy.

Cô bắc chiếc loa phóng thanh mà cô mượn được từ Subaru, hô hào tất cả mọi người, bằng cả tiếng Nhật và tiếng Etsunan: ``\gradientRGB{ Chào mừng mọi người đến với dịch vụ rút quẻ của vu nữ số một Miko-chi đây ne\~{}! Ai muốn biết vận mệnh của mình
  trong năm mới thì nhanh nhanh xếp hàng đi nào!}{0,0,0}{255,0,0}''

Rushia là người đầu tiên tiến tới, ánh mắt của cô toát ra một vẻ háo hức xen lẫn tò mò:
``Thật không? Rút một quẻ là biết ngay tương lai á? Thử xem nào!''

Miko cười rạng rỡ, chìa hộp quẻ ra: ``Rút đi, Rushia-chan! Nhớ lắc đều trước nha\~{}''
Nàng Pháp sư lắc mạnh hộp quẻ, những tiếng lách cách vang lên, càng lúc càng lớn. Cuối
cùng, cô rút ra một mảnh giấy nhỏ và mở ra đọc.

``Ể? `Tiểu hung (小凶)'? Nè!! Là sao chứ!? Không công bằng chút nào!'' nàng Pháp sư
tỏ vẻ bất bình với kết quả của tấm quẻ mà cô bốc được.

Miko hơi bất ngờ với phản ứng mãnh liệt của cô ấy. Cô đánh mắt sang phải, tìm một lý
do chính đáng để hạ hỏa đối phương lại: ``À thì\dots\ nghĩa là năm nay có thể gặp chút rắc
rối nhỏ, nhưng mà yên tâm, chỉ cần cố gắng thì mọi chuyện sẽ ổn cả thôi ne\~{}!''

``Ý của chị là em sẽ gặp xui đúng không!?'' Rushia cau mày, sồn sồn bước tới.

Ngay lúc này, Noel bước tới quầy của nàng Vu nữ, cứu nguy cho chủ quầy: ``Để tôi thử
xem! Biết đâu vận may của tôi lại tốt hơn.''

Nàng Hiệp sĩ rút quẻ một cách cẩn thận. Cô mở mảnh giấy trên đó ra và đọc to: ``Ể\dots\ `Trung cát (中吉)\dots'? Hình như không tốt mà cũng chẳng xấu nhỉ?''

``Đúng rồi đó! Trung cát nghĩa là cứ điềm tĩnh, mọi thứ sẽ dần tốt lên. Chị hợp với sự ổn
định á, ne\~{}''

Noel gật gù, hài lòng với lời giải thích của cô nàng, trong khi Rushia thì mặt mày vẫn cau
có, không phục với kết quả cô bốc.

Aki tiến tới, vẻ mặt tràn đầy tự tin: ``Tới lượt em nhé! Chắc chắn là đại cát rồi!''
Nàng Dị giới nhanh chóng rút quẻ, mở ra hí hửng. Nhưng rồi, cô nhìn chằm chằm vào
tấm quẻ mà cô đang cầm trên tay.

``Ơ\dots\ Chữ gì đây? Chữ nghĩa gì thế này?'' nàng Dị giới xoay xoay chiếc quẻ, nghiền mắt
cố đọc chữ trên đó, nhưng rồi cũng đành lắc đầu bó tay.

Miko nhìn nàng Dị giới hoang mang, bật cười: ``Ahaha\~{} Đây là chữ Hán Nôm đó Aki-chan!
Để tôi dịch cho.''

Nàng Vu nữ nhìn liếc qua chiếc quẻ của Aki, rồi nói trong sự tươi vui: ``Hmm\dots\ Của em
là `nhiều may mắn,' ứng với `đại cát (大吉)!' Xin chúc mừng\~{}!''

Aki vui mừng ôm lấy Miko, trong khi Rushia thì phồng má bất mãn: ``Tại sao mọi người
ai cũng gặp may hết vậy!? Hồi đầu năm tôi cũng đã bị cho `tiểu tung' rồi đó-nanodesu!''

Sora dịu dàng bước tới với nụ cười nhẹ. Cô tiến tới quầy của Miko, xóc lấy hộp quẻ một
cách nhẹ nhàng. ``Tôi cũng muốn thử xem lần này có gì đặc biệt không đây?''

Một chiếc quẻ lòi ra. Miko cầm lấy chiếc quẻ và đọc to mảnh giấy trên đó: ``Đại cát nữa!
Sora-chan quả nhiên lúc nào cũng may mắn hết ne\~{}!''

Nàng Thần tượng mỉm cười, ôm chiếc quẻ vào lòng, vui vẻ đáp: ``Cảm ơn cậu, Miko-chan.
Chúc cậu cũng có một năm may mắn nhé!''
Tiếp theo là Kuon và A-chan, cả hai cùng bước tới quầy để thử vận may.

\begin{dialogue}
  \speak{Kuon} Chắc cũng nên thử xem sao, biết đâu lại có điều thú vị. Đầu năm em đã
  không rút quẻ rồi.
  \speak{A-chan} Không biết Sakura-san sẽ ban ơn gì cho chúng ta đây.
\end{dialogue}

Nàng Đeo kính và cậu Tổng lần lượt rút một quẻ trong chiếc hộp. Cậu con trai mở mảnh
giấy trong chiếc quẻ của mình ra, để Miko đọc hai chữ trên đó: ``Kuon-senpai là\dots\ `Tiểu
cát (小吉)!' Năm nay sẽ là năm suôn sẻ với anh đấy.''

Cậu Tổng cầm lấy mảnh giấy trên đó, khẽ gật đầu, mỉm cười: ``Thế cũng được. Cảm ơn
em, Miko-chan.''

A-chan cũng nhanh chóng rút một quẻ, đọc qua và mỉm cười: ``Của tôi là `Trung cát
(中吉)' à. Như vậy là ổn rồi.''

Rushia nhìn hai người họ bốc được quẻ may mắn, càng tỏ vẻ bất mãn hơn. ``Vậy chị giỏi
thì chị bốc một quẻ đi, Miko-chan!'' nàng Pháp sư tỏ vẻ không phục.

``Được thôi! Nói trước, chị đây là một thần tượng vu nữ ưu tú đấy, ne!'' Miko chấp nhận
lời thách thức của cô ấy, lắc hộp quẻ một cách mạnh mẽ.

``Để xem vu nữ ưu tú này năm nay có may mắn không đây\~{}'' cô lắc thêm một hồi nữa,
rút ra một quẻ. Cô hí hửng cầm tờ giấy trên đó lên, chắc mẩm: ``\textit{Năm nay chắc là sẽ may
  lắm đây\dots}''

Nhưng, hai chữ trên mảnh giấy từ chiếc quẻ cô nhận được lại khiến cô nàng cảm thấy như
sét đánh ngang tai. Nụ cười trên môi cô dần méo xệch đi, thay vào đó là một ``cơn mưa''
mồ hôi hột.

Miko vội vàng giấu mảnh giấy sau lưng, cố làm ra vẻ bình thản. Tất nhiên, vẻ mặt lo lắng
của cô không thể qua mặt được mọi người.

\begin{dialogue}
  \speak{Aki} Chị được gì vậy?
  \speak{Kuon} Sao em lại tỏ vẻ lo lắng thế?
  \speak{Noel} Thôi nào, nói cho bọn này đi mà!
\end{dialogue}

Rushia thậm chí còn bỡn cợt: ``Chắc chị ra được quẻ nào đó đen đủi lắm nên mới giấu
bọn này chứ gì?'' Miko tỏ vẻ quả quyết: ``K-Không có đâu ne! Chỉ là chị\dots\ không muốn
nói thôi!''

Không để nàng Vu nữ phản ứng, nàng Pháp sư đã giật phắt tấm giấy mà cô nàng giấu ở
sau lưng. Cô đọc to hai chữ trên đó lên, giọng tỏ vẻ sảng khoái: ``Miko-senpai là `Đại
hung (大凶)'-nanodesu!''

Cả nhóm phá lên cười, trong khi nàng Vu nữ mặt đỏ bừng lên vì xấu hổ.

``Ha! Em tưởng chị hơn em thế nào!'' nàng Pháp sư vênh mặt lên, châm chọc ngược lại
tiền bối mình.

Miko nghe vậy cũng chỉ có thể che mặt mình vì xấu hổ, không ngờ rằng lại bị chơi khăm
đến vậy. ``\textit{Thần ơi, con đã làm gì sai đâu chứ\dots}''

\ct

Trong lúc các thành viên khác bận rộn khám phá các quầy hàng Tết, Shion và Aqua lại bị
thu hút bởi một gian hàng nhỏ với biển hiệu: ``\uline{Nặn tò he - Nghệ thuật truyền thống.}''

Cả hai nhìn chăm chú vào người nghệ nhân đang khéo léo nặn từng chi tiết trên những
con tò he đủ màu sắc, từ chú trâu vàng, hoa mai đến cả những nhân vật cổ tích.

``Tò he? Là cái gì vậy nhỉ?'' nàng Hầu gái ngước nhìn thứ mà người nghệ nhân đang bóp
nặn trên một cây que.

``Tôi cũng không biết nữa. Nhưng mà trông cũng đẹp lắm á,'' nàng Phù thủy nói, gật gù
với thứ mà người nghệ nhân ấy đang chế tạo.

``Bà nghĩ chúng ta có thể thử được không?''

``Ừ, được chứ. Nhìn nó có vẻ cũng dễ mà.''

Aqua gật đầu nhiệt tình. Cả hai nhanh chóng ngồi xuống bàn, nơi người nghệ nhân mỉm
cười chào đón họ.

``{\color{red} Các cháu muốn nặn gì? Bác sẽ hướng dẫn từng bước,}'' người đàn ông tầm tuổi trung
niên nở một nụ cười hiền từ, đưa cho hai cô gái hai cây gậy.

Aqua nói với sự năng nổ: ``C{\color{red} háu muốn làm một nhân vật giống cháu được không ạ?}''

Người nghệ nhân nở một nụ cười nhẹ: ``Được chứ! Cháu cứ làm từ từ, không cần phải
vội đâu.'' Ông hướng sang Shion, hiền từ hỏi: ``{\color{red} Còn cháu tóc trắng này? Cháu muốn làm
  con gì nào?}''

Nàng Phù thủy ngẫm nghĩ một lúc lâu, rồi bất chợt nở một nụ cười gian: ``{\color{red} Cháu muốn
  làm\dots\ một con vịt ạ! Để cho bạn cháu đấy mà.}''

Người nghệ nhân mỉm cười, đưa cho hai cô gái những bột gạo có nhiều màu sắc. Ngay
sau đó, cả hai bắt đầu nhào nặn cục bột gạo đã được cho đó một cách tỉ mỉ và cần mẫn.

Aqua cẩn thận tạo hình tóc xoáy và trang phục hầu gái của mình, trong khi Shion thì chỉnh
trang lại chiếc mỏ của con vịt trắng đội mũ lưỡi trai, mang dáng dấp của Subaru.

\dots Mặc dù, cái mỏ vịt mà cô nàng đang nặn trông chẳng khác gì một chiếc lá dính bừa
hơn.

``Hành tím'' nhìn ``Tỏi'' đang chật vật nặn cái mỏ vịt kia, liền khẽ bật cười hỏi: ``Shionchan,
bà đang nặn cái gì thế?''

Bạn của cô nhướng mày, tỏ vẻ quả quyết: ``Hả? Thì đang nặn con vịt đây chứ còn gì?
Không thấy cái mỏ này à?''

Nhìn thấy chiếc mỏ vịt khá méo mó, nàng Hầu gái che miệng lại, tủm tỉm cười: ``Sao tôi
lại thấy nó trông giống một cục kẹo bị méo mó thế nhỉ?''

Nàng Phù thủy nghe cô bạn móc máy như vậy, liền phồng má: ``Kẹo méo mó cái nỗi gì!
Bà thử làm đi rồi biết nó khó cỡ nào!''

\ct

Dần dần, Aqua cũng gặp khó khăn với đôi tay lóng ngóng. Khi tạo đến đôi mắt, cô làm
quá to khiến khuôn mặt nhân vật trông kỳ quặc.

``\dots Sao nhìn như\dots\ nhân vật này sắp khóc vậy?'' cô nàng nhăn mặt nhìn cặp mắt con tò
he của cô trông ``lạc quẻ'' so với thân hình nhỏ bé.

Shion được vậy, cười lớn châm chọc ngược lại: ``Vừa nãy bà còn chê cái của tôi xấu mà?
Cái của bà trông nó còn hài nữa!''

Nàng Hầu gái đỏ mặt, che đi sản phẩm mình đang làm: ``C-Cấm cười!''

Người nghệ nhân nhìn hai cô nàng đang nô đùa nhau ở bên cạnh chỉ cười nhẹ. Ông sau
đó cũng đã giúp họ sửa một vài chi tiết bị hỏng.

Cuối cùng, sau một hồi vật lộn, cả hai cũng hoàn thành tác phẩm của mình. Aqua thì nặn
được một ``phiên bản mini'' với vẻ mặt khá ngộ nghĩnh, còn Shion thì nặn ra được một
con vịt cũng mang dáng dấp của Subaru - mặc dù trông nó hơi khó nhìn, đến độ người ta
phải mất một lúc mới đoán ra.

Nàng Hầu gái nhìn sản phẩm mà bạn cô nặn ra, cười khì: ``Thế nào Shion-chan? Tác
phẩm của tôi trông dễ thương chứ?''

Nàng Phù thủy cười mỉm, giả vờ vỗ tay như thể không quan tâm: ``Rồi, rồi, bà là nhất rồi.
Của tôi xấu um, được chưa?''

``Nhưng mà cái của bà có thể dùng để trêu Subaru-chan được mà nhỉ?'' Aqua nhoẻn miệng
cười, giọng pha một chút tinh nghịch.

``Chứ sao! Subaru sẽ rất vui khi được tặng món quà này đấy!'' Shion che miệng cười đầy
ẩn ý, dường như cô nàng có trò để trêu chọc người bạn cùng thế hệ.

Cả hai bật cười với sản phẩm họ có, còn người nghệ nhân mỉm cười, nói: ``{\color{red} Quan trọng
  hai cháu thấy vui là được. Hai cháu làm giỏi lắm!}''

Sau khi hoàn thành, cả Subaru và Aqua ôm thành phẩm của mình đi khoe khắp nơi, khiến
mọi người không khỏi bật cười trước những ``kiệt tác'' có một không hai.

Tất nhiên là trừ Subaru, khi Shion châm chọc cô bằng tác phẩm của mình, như một thói
quen, nàng Tomboy thốt lên: ``\textbf{Đã bảo rồi, tôi không phải là vịt!!}''

\ct

Matsuri không phải là người duy nhất muốn trổ tài làm ông đồ. Trong một góc khác gần
Choco và Matsuri, một gian hàng ``Thư pháp ngày Xuân'' viết bằng chữ Quốc ngữ cũng
thu hút không ít người tham gia.

Nhìn những nét chữ được viết nắn nót và đẹp đẽ trên những tấm giấy đỏ, nàng Song tính
nghĩ thầm: ``\textit{Viết chữ đẹp như thế này, chắc mình làm được!}''

Không do dự, cô nàng tiến đến bàn, nơi một ông đồ đang kiên nhẫn viết chữ cho khách.
Ngay sau khi vị khách rời đi, cô nàng Trung học hớn hở bước đến gần. Ông đồ mỉm cười
một cách thân thiện: ``{\color{red} Cháu gái muốn viết chữ gì để treo ngày Tết nào?}''

Haato tiến tới, tay chỉ vào những tờ giấy đỏ, hớn hở đáp: ``{\color{red} Cháu không đến đây để xin
  chữ. Cháu muốn tự viết được không ạ? Cháu muốn thử tay nghề một chút!}''

Ông đồ khẽ ngạc nhiên với yêu cầu khác thường từ cô nàng Trung học đó, nhưng rồi mỉm
cười, đưa cô một tờ giấy đỏ và một cây bút lông. ``Được chứ. Cháu muốn viết chữ gì, bác
sẽ hướng dẫn từng nét.''

Cô nàng ngẫm nghĩ một lác, sau đó tươi cười nói: ``{\color{red} Chữ `Phúc (福)' ạ! Cháu thấy chữ đó
  đẹp và ý nghĩa lắm ạ.}''

Ông đồ gật đầu, sau đó chỉ cô cách cầm bút lông đúng cách. Haato chăm chú quan sát và
bắt đầu thử viết.

``\textit{Chắc cũng giống như viết Kanji bên mình thôi! Cái này dễ ẹt à!}'' nàng Song tính thầm
nghĩ, cười khoái chí khi nhìn những nét nắn nót của ông đồ, tin rằng cô cũng có thể viết
được như ông ấy dựa vào trình độ viết Kanji của mình.

Nhưng, đời không như là mơ với nàng Trung học này.

Ngay khi Haato cầm lấy cây bút lông, cô bắt đầu loay hoay, như có ý chống lại sự khéo
tay của cô nàng. ``Sao cái bút này khó điều khiển vậy!? Nét cứ lệch qua lệch lại thế này!''

Những nét chữ trên giấy đỏ mà cô nàng đang viết trông không những không giống chữ
``Phúc,'' mà nó còn giống như\dots\ một con rắn đang uốn éo.

Cô nàng bắt đầu tỏ vẻ bối rối, trái ngược hẳn với sự tự tin ban đầu của mình: ``{\color{red} Ê-tồ\dots\ Cháu nghĩ làm mình đã làm sai chỗ nào ạ\dots?}''

Ông đồ nhìn Haato đang ngơ ngác ấy, khẽ bật cười: ``{\color{red} Nhìn cách cháu viết thế này chắc
  cháu cũng thạo chữ Hán lắm nhỉ? Nhưng mà, viết thư pháp thì không giống như viết trên
  vở đâu. Quan trọng là cháu phải có tấm lòng và sự kiên nhẫn cơ.}''

``{\color{red} À\dots\ ra thế.}'' Cô nàng hít sâu, quyết tâm hơn. Lần này, cô nàng viết lại một chữ `Phúc'
khác cẩn trọng hơn.

Nhưng lần này, do sự cẩn trọng quá mức mà cô nàng lại quá căng thẳng, khi tay cô cầm
cây bút một cách run rẩy để lại những vệt mực không đều.

Nàng Song tính lúc này đang nghiến răng, tay cố viết thật nắn nót những nét cuối cùng
với tràn đầy sự quyết tâm. Trong khi đó, ông đồ chỉ có thể cười trừ nhìn cô nàng ấy.

``Xong rồi!'' cô nàng thốt lên, giơ hai tay lên trời, tò mò xem thành phẩm của mình. Nhưng
rồi\dots

``{\color{red} Ể!? Tại sao nét cuối lại to thế này!?}'' cô thốt lên với sự ngỡ ngàng khi chữ cô viết lại
không được như ý. Haato ngước mắt nhìn ông đồ, mắt rưng rưng.

``{\color{red} Bác ơi, cháu lại làm hỏng mất rồi\dots}''

Ông đồ nhìn chữ viết của cô nàng, mỉm cười một cách dịu dàng, rồi nói với sự thông cảm:
``{\color{red} Chữ không cần phải hoàn hảo, cháu à.}''

``{\color{red} Thật ạ\dots?}''

``Ừ. Quan trọng là cháu đã dồn tâm sức vào. Đây chính là cái hồn của thư pháp đấy.''

Nghe xong những lời động viên từ ông đồ, cô nàng cảm thấy an ủi phần nào. Cô nhìn vào
cặp mắt của ông ấy, tràn đầy sự quyết tâm. ``{\color{red} Được rồi, cháu sẽ thử viết lại lần nữa!}''

\ct

Cuối cùng, cô nàng nỗ lực thêm một lần nữa, và thành công viết được chữ ``Phúc'' đơn
giản nhưng rõ ràng. Ngay khi thành phẩm của mình trông ưa nhìn hơn, Haato cầm giấy
đỏ lên trưng cho mọi người, như thể cô vừa tìm thấy một loại kho báu quý giá: ``{\color{red} Cháu
  làm được rồi! Tuy không đẹp lắm, nhưng cháu rất tự hào ạ-chama!}''

Ông đồ nhìn cô, gật đầu hài lòng. ``{\color{red} Vậy sao cháu không đi khoe cho các bạn cháu nhỉ?}''

Dứt lời, cô nàng cầm tờ giấy đỏ đó rời khỏi gian hàng, không quên cảm ơn ông ấy.

``\textbf{Mọi người ơi, nhìn này! Đây là chữ do chính Haachama ta đây viết này!!}''

Cầm chữ ``Phúc'' trên tay, Haato hào hứng vui vẻ đem đi khoe với các thành viên khác.
Đa phần đều cảm thấy ngưỡng mộ và khen ngợi tác phẩm do chính cô tạo ra.

Ngoại trừ một vài người thích châm chọc cô nàng Trung học này. Nàng Hầu gái nhìn chữ,
nghiêng đầu khó hiểu: ``Ủa? Sao trông giống cái mặt cười vậy?''

Nhận lời bình từ nàng Hầu gái kia, Haato phồng má lên, phụng phịu: ``Mặt cười gì chứ!
Đây là chữ `Phúc'!''

Kuon nhìn tác phẩm của cô nàng, cũng bật cười trêu: ``Phúc kiểu này là phúc `trườn bò'
rồi chứ có phải là phúc may đâu!''

Nàng Song tính hét lớn: ``\textbf{Em không có vẽ con rắn!! Hai người đúng là quá đáng!}''

Tiếng cười của mọi người rộn ràng vang lên, trong khi Haato thì rơm rớm nước mắt vì lời
trêu chọc của hai người.

\il

Ngày Tết tại đền Vương, với không khí trang trọng nhưng cũng đầy ắp tiếng cười, dần trôi
qua trong những hoạt động vui tươi và mang đậm bản sắc dân tộc. Các cô gái từ \hll\,
từ những hoạt động thú vị đến những trải nghiệm mới lạ, tất cả đều cảm thấy gắn bó và
hứng khởi với không khí này.

Miko và Matsuri bắt đầu tất tả dọn đồ, chuẩn bị về nhà của Kuon để ăn bữa tối Tết cuối
cùng. Nàng Vu nữ vừa cắp nách chiếc quẻ, vừa nhìn mọi người: ``Cảm ơn mọi người đã
tham gia. Chị nghĩ hôm nay thật sự là một ngày Tết đáng nhớ, ne!''

Kuon cười nhẹ, đáp: ``Không chỉ có mình các em đâu, mà anh nghĩ cả các vị khách đến
đền hôm nay cũng rất vui vẻ. Nhắt là một số người cảm thấy bất ngờ trước sự hiện diện
của mấy đứa đấy.''

A-chan ở bên cạnh cậu đang phụ giúp mọi người dọn dẹp, cũng mỉm cười: ``Đúng là Tết
là dịp để đoàn tụ và thể hiện sự gắn kết nhỉ.''

Những tiếng cười đùa vang lên khi các thành viên đi qua các gian hàng, mỗi người đều
bày tỏ sự thích thú với hoạt động mà mình tham gia. Lúc này, có vẻ như không khí ngày
Tết đã len lỏi vào từng người, từ những cái nhìn tươi vui đến những câu chuyện đầy hào
hứng, dù rằng những ngày Tết giờ đây đang trôi dần về những khoảng thời gian cuối cùng.

Matsuri nhìn về phía mọi người, hào hứng nói: ``Có lẽ năm sau chúng ta sẽ quay lại nơi
này. Mọi người chúng ta ai nấy đều đã có kỷ niệm đẹp về những ngày Tết Nguyên Đán
đã qua chứ?''

Sora gật đầu, chắp tay cầu chúc mọi người. ``Đúng thế. Chị chắc chắn sẽ nhớ mãi màn
múa lân, màn bắn pháo hoa rực trời và hành trình về quê nữa.''

Rushia, dù vẫn còn hơi hậm hực về cái chuyện bốc quẻ, nhưng cô vẫn thỏa mãn: ``Cái quẻ
của tôi có vẻ không ổn lắm, nhưng ít ra nó khiến tôi thấy may mắn hơn ai đó đằng kia.''

``Đừng có mà khịa chị chứ\dots'' Miko nhăn mặt, tỏ vẻ không vừa ý khi cô nàng nhắc lại
chuyện đáng xấu hổ đó.

Mọi người cùng nhau đi dạo một vòng quanh đền, tay cầm những món đồ nhỏ như quạt
giấy, tranh Tết, và những bức thư pháp vừa viết xong. Không gian tĩnh lặng của đền, kết
hợp với những tiếng pháo và tiếng cười từ nhóm bạn, tạo nên một bức tranh thật sự tuyệt
vời về một ngày đầu năm.

Aki ngắm nghía lại bức câu đối mà cô được một ông đồ viết tặng trong tay, cười nhẹ:
``Đây thực sự là một chuyến đi đáng nhớ. Mình hy vọng sẽ có nhiều dịp như vậy để mọi
người cùng nhau trải nghiệm.''

Fubuki hớn hở đáp: ``Ngày Tết mà, chỉ cần có mọi người bên cạnh thì mọi thứ đều vui vẻ
hết.'' Mio nói chen vào: ``Miễn là đừng làm loạn tới mọi người là được.''

Và thế là, sau một ngày đầy ắp những hoạt động mang đậm màu sắc Tết, các thành viên
\hll\ trở về với những niềm vui và sự hào hứng trong lòng. Dù có chút vất vả, có lúc
bối rối vì chưa quen với những tập tục, nhưng tất cả đều cảm nhận được một điều: cái
Tết, dù có xa lạ đến mấy, cũng luôn mang đến sự ấm áp, tình thân và niềm vui không thể nào quên.

\ct

\begin{center}
  \uline{7:00 tối\\Tầng ba - Nhà Hikawa.}
\end{center}

Mùng Ba Tết, ngày cuối cùng mà các thành viên \hll\ sẽ hoạt động trải nghiệm Tết tại
đây, mọi người quây quần bên bàn ăn, thưởng thức một bữa tiệc giao thoa giữa ẩm thực
cổ truyền của Etsunan và Nhật Bản.

Kaien bày ra trước mắt những chiếc bánh chưng cuối cùng mà mọi người (cả \hll\ và
bên ngoại) đều đã dồn hết công sức gói trước bàn ăn. ``Cảm ơn các cháu đã bỏ công sức
chuẩn bị. Bữa tối hôm nay thật đẹp mắt.''

Ryuuhou ngồi bên cạnh mẹ, hớn hở nói: ``Sora-onee-chan và Ayame-onee-chan lần này
đã trổ tài làm món nem rán này đấy ạ! Nhưng hình như Ayame-onee-chan ăn mất nửa số
nem rồi\dots''

Nàng Quỷ sừng, vẫn còn dính sợi miến và mảnh bánh tráng trên miệng, giả vờ ngơ ngác:
``Ơ, chị có ăn đâu\dots\ Mà chắc chỉ nếm thử một chút thôi!''

Kuon nhìn mép của cô nàng còn dính đồ ăn, cười trừ: ``Mồm mép em dính đầy thế kia
còn giấu được ai.''

Sora cười mỉm trước nàng Oni háu ăn đó: ``Em không ngờ làm nem lại khó đến vậy.
Nhưng nhờ có Kaien-san chỉ dẫn mà em đã học được cách cuốn rồi.''

Ngoài những món ăn truyền thống xứ Etsu, những món ăn Nhật Bản truyền thống khác
cũng đươc bài trí một cách khéo léo.

Korone cầm đĩa sushi lên, hào hứng khoe trước mặt mọi người: ``Đây là tác phẩm của em
với Okayu! Bọn em đã cố gắng làm đẹp nhất rồi đó ha\~{}''

Okayu cũng gật gù tự hào, khuôn mặt hiện lên vẻ tự mãn: ``Có cả cá hồi mà Kuon-senpai
đã mua giúp nữa đấy. Ngon lắm luôn!''

Cậu Tổng nhìn đĩa sushi đó, gật gù: ``Phải công nhận là hai đứa làm rất tốt,'' đồng thời
không quên bông đùa: ``Thế này thì khách sạn cũng phải mời về làm bếp trưởng thôi.''

Tất cả mọi người bật cười với câu nói của cậu chàng. Ở một diễn biến khác, Matsuri ngắm
nghía những đĩa thức ăn đầy thịnh soạn: ``Chị thì làm món gì cũng không đẹp được như
mấy đứa\dots\ Nhưng mà hương vị chị làm thì ngon lắm nha!''

Aqua nghe cô nàng nói vậy, khẽ hoảng hốt: ``Ơ, Matsuri-senpai\dots\ đừng nói món súp rong
biển của chị là\dots\ món em vừa nếm thử lúc nãy chứ!?''

Nàng Cổ động chỉ nháy mắt: ``Đoán xem\~{}'', càng khiến cho nàng Hầu gái toát mồ hôi hột
vì lo sợ hơn.

Ở đầu bàn, những món ăn cổ truyền Etsunan khác như bánh chưng, thịt đông, canh măng
và dưa hành cũng khiến các cô gái không khỏi ngạc nhiên.

Mio nhìn đĩa bánh chưng mà mẹ của Kuon bưng ra, khen ngợi: ``Bánh chưng này ngon
thật! Nhưng cứng quá, phải cắt làm đôi mới ăn được. Ăn bao nhiêu lần cũng không thấy
chán!''

Mel ngồi bên cạnh cô, mỉm cười: ``Chị lại thích vị béo ngậy của nhân thịt á. Phần nếp
thì\dots\ cũng tạm ổn thôi.''

Noel cầm miếng thịt đông, đôi mắt sáng lóa lên: ``Món này trông như thạch vậy! Nhìn nó
núng nính thế này trông đã mắt ghê!''

Cậu con trai nghe nàng Hiệp sĩ khen lấy khen để món thịt đông do bố cậu làm, cười trừ:
``Anh không thích ăn thịt đông lắm, nên anh cũng không hiểu em đang nói gì\dots\ Nhưng
mà miễn em thấy ngon là được.''

Shion liếc nhìn món thịt đông kia, tỏ vẻ thắc mắc: ``Không biết chị có thể làm được món
này bằng phép thuật của mình được không nhỉ\dots\ Mà thôi, chắc là không cần đâu.''

Rushia ngồi đối diện cô, lắc đầu cười: ``Em nghĩ là không đâu. Nhưng mà Shion-senpai
có thể làm thử món bánh chưng mini mà Kuon-senpai đã chỉ hôm nay xem.''

``Ý hay đấy.''

Ở giữa bàn, không khí càng rộn ràng hơn khi các cô gái chia sẻ về ngày Tết cuối cùng.

Fubuki cười lớn, ôn lại những ngày náo nhiệt mà cô cùng họ đã trải qua: ``Những ngày
vừa rồi thật sự là một ngày đáng nhớ. Từ những món ăn đặc trưng cho đến các hoạt động
ngày Tết\dots\ Tất cả đều thú vị á!''

Haato - một trong những người hăng hái nhất trong những ngày vừa rồi - cũng tươi cười:
``Nhất là phần rút quẻ! Tôi vẫn còn giữ quẻ đây này! Tôi vẫn nhớ cái hôm mà chúng ta
đã đua xe trên quốc lộ ngày mùng Một à!''

Miko cười lớn: ``Chị vẫn nhớ ngày mùng Hai khi mà chúng ta chơi Hanetsuki à! Mấy
đứa đúng là đánh ác liệt thật!''

Aki mỉm cười: ``Đúng thật đó. Năm ngày Tết trôi qua nhanh thật đấy.''

Ryujin ngồi ở phía đầu bàn, nhìn các cô gái và gia đình mình với ánh mắt đầy tự hào. Ông
nâng cốc bia lên, nói với tất cả mọi người: ``Nhân dịp hôm nay cũng là ngày cuối cùng
các cháu tới đây, chú muốn gửi lời cảm ơn mọi người vì đã mang lại không khí Tết cho
gia đình Hikawa nhà chú.''

Kaien tiếp lời: ``Cô cũng muốn cảm ơn các cháu vì đã tới đây chung vui với cô chú trong
những ngày vừa qua.''

Kế tiếp đó là Ryuuhou: ``Dù năm ngày vừa qua là khoảng thời gian ngắn ngủi, nhưng em
mong các chị đã học được nhiều điều về lễ hội Tết Nguyên Đán tại Etsunan này!''

Cuối cùng là Kuon, kết thúc bài cảm ơn: ``Anh chúc tất cả mọi người ở \hll\ một năm
mới an khang thịnh vượng, năm mới có nhiều lộc, nhiều may mắn trong cuộc sống!''

Tất cả mọi người đồng loạt vỗ tay, với nàng Đeo kính đang cầm chiếc máy quay ghi lại
những thước phim cuối cùng cho bộ phim Tết của họ sắp tới.

``Thật tuyệt vời khi được tham gia vào những trải nghiệm thế này. Tôi nghĩ ai cũng sẽ nhớ
mãi về dịp này,'' nàng Đeo kính nói, nở một nụ cười tươi rói khi tất cả hai mươi tư người
họ đều vui vẻ như vậy.

``Cảm ơn mọi người vì đã là một phần trong chuyến đi đầy ý nghĩa này!'' cậu con trai chốt
hạ lại, theo sau là một tràng pháo tay khác giòn giã hơn.

Kết thúc bữa ăn, mọi người cùng nhau chụp một bức ảnh kỷ niệm. Dưới ánh đèn trắng
ấm áp, những cành hoa đào và những cành quất đã chín ươm, các thành viên \hll\ và
gia đình Hikawa mỉm cười thật tươi, khép lại một ngày Tết trọn vẹn trong tiếng cười và
niềm vui lan tỏa.

\il

Sau bốn ngày đầy ắp tiếng cười và trải nghiệm mới mẻ, cả nhóm đã thu về những thước
phim tư liệu quý giá. Những cảnh quay từ các hoạt động tại nhà Hikawa, chuyến đi đền
Vương, những buổi dọn dẹp nhà cửa mệt nhọc nhưng vui vẻ, những chuyến đi hỗn loạn
nhưng sôi nổi, cho đến những màn trình diễn đậm chất \textit{hologra} đều được lưu giữ cẩn thận,
hứa hẹn tạo nên một bộ phim đặc biệt dành tặng khán giả.

\begin{center}
  \uline{Ngày hôm sau.}
\end{center}

Sáng sớm hôm sau, khi ánh mặt trời vừa ló dạng trên bầu trời thành phố của Etsunan, các
cô gái gói ghém hành lý, chuẩn bị trở về thành phố Điện tử để tiếp tục công việc thường
nhật.

Đứng trước cánh cửa thần kỳ mà họ đã dùng để tới đây, tất cả mọi người tay xách nách
mang đồ ăn mà họ đã mang tới, kèm theo đó là những món quà quê từ gia đình Hikawa.

Đích thân Kuon là người đứng tiễn đoàn người đó. Khuôn mặt cậu hiện ra một sự vui vẻ
mãn nguyện, phần vì các cô gái đã hoàn thành xuất sắc nhiệm vụ, phần vì\dots\ họ cuối cùng
cũng trả lại không gian bình yên cho căn nhà.

``Các em đi đường cẩn thận nhé. Anh sẽ ở lại vài ngày để hoàn thành báo cáo.''

Sora mỉm cười dịu dàng, khẽ gật đầu: ``Cảm ơn anh và gia đình đã đón tiếp chúng em
nồng nhiệt như vậy. Em rất vui vì có được một trải nghiệm đáng nhớ như thế.''

Mio vẫy tay, vừa là để hối thúc quay về, vừa nói lời chào tạm biệt tới cậu Tổng: ``Chúng
em sẽ đợi xem thành quả từ bộ phim này! Hẹn gặp lại anh ở công ty!''

Tất cả mọi người lần lượt đi qua cánh cửa thần kỳ, không quên gửi lời cảm ơn đến bố mẹ
và em gái Kuon vì đã dành cho họ một cái Tết ấm áp, ý nghĩa.

Nhìn cánh cửa thần kỳ dần biến mất, không khí xung quanh trong căn nhà gia đình Hikawa
cũng trầm mặc hẳn. Sự im lặng yên bình vốn có đã thay thế sự bát nháo trong những ngày
vừa qua.

Cô em gái của Kuon nhìn cánh cửa dần biến mất, khuôn mặt cô tỏ vẻ một chút tiếc nuối,
nhưng vẫn tỏ ra vui vẻ: ``Những cô gái ấy thật đặc biệt, đúng không bố?''

Ryujin gật gù: ``Ừ. Bố hy vọng họ có thể sẽ trở lại với nhiều câu chuyện mới.''

Kaien mỉm cười: ``Có khi sau này một trong số họ sẽ làm bạn gái của con thì sao! Đúng
không, Kuon!''

Kuon đỏ mặt, ngượng ngùng đáp: ``Mẹ đừng có trêu con nữa đi\dots\ Con chưa sằng kết hôn
đâu ạ\dots''

Cả gia đình Hikawa quay trở về với nếp sống bình thường, tận hưởng nốt những ngày Tết
bình yên hiếm hoi. Tuy có chút tiếc nuối khi chia tay mọi người, nhưng cậu cũng cảm
thấy nhẹ nhõm vì mọi thứ đã diễn ra tốt đẹp.

Đúng là một ngày Tết hỗn loạn nhưng cũng đầy ý nghĩa!

\end{document}